% =====================================================================
%  KG-Commit --- Experimental Setup (DETAILED DRAFT for the Q1 paper)
%
%  This is a deliberately verbose, fully-explained draft covering the FINAL
%  (V4) methodology of the project. It is meant to be read and then distilled
%  into a tighter camera-ready "Experimental Setup" section. Nothing here is
%  length-constrained; every design choice is spelled out with its rationale
%  and the exact constants used, so the summarised version can be both accurate
%  and defensible under Q1 review.
%
%  STANDALONE BUILD: compiles on its own with `pdflatex experimental_setup_draft.tex`
%  (run *twice* to resolve cross-references).
%
%  TO \input THIS INTO THE MAIN DOCUMENT INSTEAD (e.g. KG-Commit evolution.tex,
%  which already defines \code/\node/\edge and loads booktabs/enumitem etc.):
%  delete everything from \documentclass down to \begin{document} below, and
	%  delete the closing \end{document} at the end of the file, then
%  % =====================================================================
%  KG-Commit --- Experimental Setup (DETAILED DRAFT for the Q1 paper)
%
%  This is a deliberately verbose, fully-explained draft covering the FINAL
%  (V4) methodology of the project. It is meant to be read and then distilled
%  into a tighter camera-ready "Experimental Setup" section. Nothing here is
%  length-constrained; every design choice is spelled out with its rationale
%  and the exact constants used, so the summarised version can be both accurate
%  and defensible under Q1 review.
%
%  STANDALONE BUILD: compiles on its own with `pdflatex experimental_setup_draft.tex`
%  (run *twice* to resolve cross-references).
%
%  TO \input THIS INTO THE MAIN DOCUMENT INSTEAD (e.g. KG-Commit evolution.tex,
%  which already defines \code/\node/\edge and loads booktabs/enumitem etc.):
%  delete everything from \documentclass down to \begin{document} below, and
	%  delete the closing \end{document} at the end of the file, then
%  % =====================================================================
%  KG-Commit --- Experimental Setup (DETAILED DRAFT for the Q1 paper)
%
%  This is a deliberately verbose, fully-explained draft covering the FINAL
%  (V4) methodology of the project. It is meant to be read and then distilled
%  into a tighter camera-ready "Experimental Setup" section. Nothing here is
%  length-constrained; every design choice is spelled out with its rationale
%  and the exact constants used, so the summarised version can be both accurate
%  and defensible under Q1 review.
%
%  STANDALONE BUILD: compiles on its own with `pdflatex experimental_setup_draft.tex`
%  (run *twice* to resolve cross-references).
%
%  TO \input THIS INTO THE MAIN DOCUMENT INSTEAD (e.g. KG-Commit evolution.tex,
%  which already defines \code/\node/\edge and loads booktabs/enumitem etc.):
%  delete everything from \documentclass down to \begin{document} below, and
	%  delete the closing \end{document} at the end of the file, then
%  % =====================================================================
%  KG-Commit --- Experimental Setup (DETAILED DRAFT for the Q1 paper)
%
%  This is a deliberately verbose, fully-explained draft covering the FINAL
%  (V4) methodology of the project. It is meant to be read and then distilled
%  into a tighter camera-ready "Experimental Setup" section. Nothing here is
%  length-constrained; every design choice is spelled out with its rationale
%  and the exact constants used, so the summarised version can be both accurate
%  and defensible under Q1 review.
%
%  STANDALONE BUILD: compiles on its own with `pdflatex experimental_setup_draft.tex`
%  (run *twice* to resolve cross-references).
%
%  TO \input THIS INTO THE MAIN DOCUMENT INSTEAD (e.g. KG-Commit evolution.tex,
%  which already defines \code/\node/\edge and loads booktabs/enumitem etc.):
%  delete everything from \documentclass down to \begin{document} below, and
	%  delete the closing \end{document} at the end of the file, then
%  \input{experimental_setup_draft.tex} from the parent .tex.
% =====================================================================
\documentclass[11pt,a4paper]{article}
\usepackage[margin=2.4cm]{geometry}
\usepackage[T1]{fontenc}
\usepackage[utf8]{inputenc}
\usepackage{lmodern}
\usepackage{amsmath,amssymb}
\usepackage{booktabs}
\usepackage{array}
\usepackage{enumitem}
\usepackage{xcolor}
\usepackage[hidelinks]{hyperref}

% ---- macros used throughout (match the main document's conventions) ----
\providecommand{\code}[1]{\texttt{#1}}
\providecommand{\node}[1]{\textsf{#1}}
\providecommand{\edge}[1]{\textsc{#1}}

% ---- helper for safe section references (avoids any parsing issues with \S\ref) ----
\newcommand{\secref}[1]{\S\ref{#1}}

\begin{document}
	
	\section{Experimental Setup}
	\label{sec:experimental-setup}
	
	This section specifies, in full, the data, the knowledge-graph construction, the
	inference models, the online evaluation protocol, the metrics, the baselines, and
	the reproducibility conditions under which every result in this paper was produced.
	Two principles govern the whole design and are worth stating up front because they
	recur throughout. \textbf{First, everything is \emph{just-in-time (JIT) and fully
			online}:} the graph grows one commit at a time in true chronological order, and a
	commit is always predicted from a model and a graph state built \emph{only} from
	strictly earlier commits --- there is no random split, no look-ahead, and no
	statistic is ever fit on data that includes or postdates the commit under test.
	\textbf{Second, the study is designed as a controlled comparison:} when we ask which
	code representation or which inference method is best, exactly one factor is varied
	at a time while the protocol, the classifier family, the warm-up, and the block
	schedule are held fixed, so any measured difference is attributable to that factor
	and not to a confound.
	
	\subsection{Subject program, commit scope, and labels}
	\label{sec:setup-data}
	
	\paragraph{Project.} All experiments are conducted on \texttt{apache/groovy}, a
	mature, long-lived Java project (the Groovy programming language and its compiler,
	runtime, and tooling). We work from a full local \texttt{git} clone, so the raw
	material for every commit --- its message, its parent(s), and the exact byte content
	of every file before and after the change --- is available offline and deterministic.
	
	\paragraph{Labelled commit scope.} The study set is the \textbf{$8{,}059$ commits}
	of \texttt{apache/groovy} that appear in the \textbf{ApacheJIT} dataset
	(\texttt{data/apachejit/projects/apache\_groovy.csv}), spanning \textbf{2003--2019}.
	ApacheJIT provides, for each commit, a \emph{buggy/benign} label derived by an
	SZZ-style procedure (a commit is buggy if a later fix commit points back to it) and
	the twelve standard JIT change metrics. These $8{,}059$ commits are the only ones we
	treat as the labelled study set and the only ones grown into the AST/structural and
	semantic layers. The surrounding history is retained in the database as
	\node{Commit} nodes for context (the full clone has $16{,}887$ commits), but only the
	ApacheJIT commits are flagged \code{in\_jit=true} and are the units of prediction.
	
	\paragraph{Class balance and its consequences.} The labelled set is imbalanced:
	\textbf{$1{,}614$ buggy / $6{,}445$ benign}, i.e.\ a $\approx\!20\%$ positive rate
	overall. Crucially, the positive rate is \emph{not stationary}: it rises from
	$\approx\!6\%$ in 2003 to a peak of $\approx\!39\%$ in 2012 and then falls back to
	$\approx\!6\%$ by 2019 (measured directly from the labels). This pronounced
	\emph{concept drift} is the single most important property of the data for our
	purposes: it is why a random or single-split evaluation is inappropriate, why we
	adopt a time-aware online protocol, and why the tail of every online trajectory
	declines (the achievable precision/recall shrink as positives become rare). We
	return to this when defining the protocol (\secref{sec:setup-protocol}) and the
	metrics (\secref{sec:setup-metrics}).
	
	\paragraph{Base snapshot.} The structural layers are seeded once at a fixed
	\emph{base} commit --- \texttt{408b29851d7bbe4d343340832297e4be7e0c5578}, the first
	substantial Java commit --- whose $120$ \texttt{.java} files (${\approx}46$k AST
	nodes) constitute the initial code structure onto which all subsequent per-commit
	changes are laid. The base is a property of the \emph{code}, not of the labelled
	set; it merely fixes the starting state of the evolving graph.
	
	\subsection{Knowledge-graph construction and online growth}
	\label{sec:setup-kg}
	
	The knowledge graph is a \emph{heterogeneous information network} (HIN): typed nodes
	and typed edges in one store (Neo4j). It fuses three families of evidence that the
	defect-prediction literature usually treats separately --- \emph{process} evidence
	(who changed what, when, and how often), \emph{code-structure} evidence (the
	syntactic/semantic shape of the change and its surroundings), and \emph{textual}
	evidence (the natural-language and lexical content of the change). The design
	rationale for a graph (rather than a flat feature table) is that anchoring each
	commit's change to \emph{shared, persistent} code entities is what makes relational
	and neighbourhood reasoning across the entire history possible.
	
	\paragraph{Core (process / relational) layer.} The Core layer is the classical
	commit graph and is present in \emph{every} experimental configuration. Its node
	types are \node{Project}, \node{Branch}, \node{Developer}, \node{Commit},
	\node{File}, and \node{Issue}; its edges encode authorship
	(\node{Commit}$\xrightarrow{\text{\edge{AUTHORED\_BY}}}$\node{Developer}), the
	history spine (\node{Commit}$\xrightarrow{\text{\edge{PARENT\_OF}}}$\node{Commit}),
	file changes
	(\edge{ADDED}/\edge{MODIFIED}/\edge{DELETED}/\edge{RENAMED\_FROM}/\edge{RENAMED\_TO}
	between \node{Commit} and \node{File}), and issue links
	(\node{Commit}$\xrightarrow{\text{\edge{FIXES\_ISSUE}}}$\node{Issue}). Every
	\node{Commit} also carries the twelve ApacheJIT change metrics (size \code{la},
	\code{ld}; dispersion \code{nf}, \code{nd}, \code{ns}, \code{ent}; history and
	experience \code{age}, \code{ndev}, \code{nuc}, \code{aexp}, \code{arexp},
	\code{asexp}) and its \code{author\_ts} (the UNIX timestamp that defines online
	arrival order), \code{buggy} label, and \code{jit\_year}.
	
	\paragraph{Structural (code) layers as \emph{delta} graphs.} On top of the Core, we
	materialise the code structure of every changed \texttt{.java} file. The key modelling
	decision --- and a central contribution --- is that a commit is not attached to a
	per-commit snapshot blob; it is attached, by four typed \emph{delta} edges
	(\edge{ADDS}, \edge{REMOVES}, \edge{UPDATES}, \edge{MOVES}), directly to the specific
	code nodes it changed, laid on top of a single evolving structure per file. Concretely,
	for each tracked file the graph holds one \emph{current} structure; when a commit
	touches the file we (i) parse the before- and after-states with the \emph{same}
	deterministic builder, (ii) diff them, (iii) write the delta edges from the
	\node{Commit} to the affected nodes, and (iv) evolve the stored structure so the next
	commit diffs against the correct state. Node identity is threaded across versions by a
	deterministic, per-file parse-local-id map, so the ``after'' map of commit $N$ is
	byte-identical to the ``before'' map of the next commit that touches the file when the
	intervening source is unchanged. Removed nodes are \emph{retained} and flagged
	(\code{alive=false}, \code{removed\_by}), so the complete change history is queryable
	while the live structure always reflects the current code. This makes storage grow
	with the \emph{amount of change}, not with history-length $\times$ file-size, and it is
	what makes the whole graph online-extensible.
	
	We instantiate this delta-graph scheme for \textbf{five} structural representations of
	increasing coarseness, all built per changed \texttt{.java} file and attached at the
	\node{File} level:
	\begin{itemize}
		\item \textbf{AST} --- the abstract syntax tree; one node per syntactic construct
		(declaration, statement, expression, operator, literal, and each identifier/operator
		leaf). Finest granularity ($136$ node types).
		\item \textbf{CFG} --- the intra-procedural control-flow graph; one node per statement
		or control construct ($18$ node types).
		\item \textbf{DFG} --- the def--use data-flow graph over local variables; one node per
		definition/use site ($7$ node types).
		\item \textbf{PDG} --- the program-dependence graph; the CFG's statement vertices plus
		control-dependence and data-dependence \emph{edges} (no new nodes; $18$ node types).
		We label it PDG rather than CPG because we deliberately keep dependence at statement
		granularity and do \emph{not} embed the AST subtree under each statement; a canonical
		code property graph would contain the AST as a subgraph and is out of scope.
		\item \textbf{Token-seq} --- a deliberately weak flat statement sequence (statement
		nodes chained by \edge{NEXT}, no control or data structure); a lower-bound rival used
		only in the representation comparison.
	\end{itemize}
	Table~\ref{tab:setup-layers} reports the materialised scale of each layer. The
	${\approx}10\times$ gap between AST ($3.70$\,M nodes) and the others ($0.26$--$0.47$\,M)
	is a pure \emph{granularity} effect: a statement such as \code{int x = foo(a,b)+1;} is
	one node in CFG/PDG but a dozen-plus nodes in the AST. CFG and PDG have identical node
	counts because the PDG adds only dependence edges to the CFG's vertex set.
	
	\begin{table}[t]
		\centering
		\caption{Materialised scale of each structural layer and its per-commit change-token
			stream (Groovy, $8{,}059$ commits). ``Token types'' counts the distinct
			\code{\{edge\}:\{node-type\}} change-tokens; ``Commits w/ tok.'' is the number of
			labelled commits that emit at least one token in that layer.}
		\label{tab:setup-layers}
		\begin{tabular}{lrrrrr}
			\toprule
			Layer & Nodes & Delta-edges & Node types & Token types & Commits w/ tok. \\
			\midrule
			AST       & 3{,}699{,}439 & 3{,}614{,}922 & 136 & 314 & 7{,}523 \\
			CFG       &   372{,}488 &   408{,}301 & 18 & 53 & 6{,}194 \\
			DFG       &   469{,}662 &   648{,}553 &  7 & 27 & 6{,}066 \\
			PDG       &   372{,}488 &   411{,}869 & 18 & 63 & 6{,}261 \\
			Token-seq &   257{,}708 &   293{,}287 & 16 & 47 & 6{,}011 \\
			\bottomrule
		\end{tabular}
	\end{table}
	
	\paragraph{The differ, and why matching is fair across layers.} The AST layer is
	diffed as a \emph{tree} by bottom-up subtree hashing (identical subtrees in the before
	and after trees are matched greedily, preferring equal source positions; unmatched
	before-nodes are deletions, unmatched after-nodes insertions, matched-but-relabelled
	nodes updates, reparented nodes moves). The four non-tree layers (CFG/DFG/PDG/Token-seq)
	are general typed graphs and are diffed by a single \emph{node-identity-lite} matcher
	applied \emph{identically} to each: matching is scoped within a method (\code{class::name},
	stable across commits) and keyed by \emph{ordinal} (the $k$-th node of a given type in a
	method maps to the $k$-th such node in the other version), which is robust to line
	shifts. Using one shared differ for the alternatives is the fairness guarantee for the
	representation comparison: only the subgraph varies, never the difference model. We
	disclose that the AST uses a tree differ while the others use the ordinal matcher --- a
	deliberate specialisation (trees vs.\ general graphs), not a hidden confound.
	
	\paragraph{Semantic-text layer (CSTG).} The third first-class layer is the
	\textbf{Commit Semantic-Text Graph (CSTG)}, which models the \emph{textual}
	description of a change (commit message $+$ diff text) as a graph grounded in code.
	Its node types are \node{Term} (a typed vocabulary term, with \code{kind} $\in$
	\{\emph{code}-identifier, \emph{bug}-lexicon, \emph{action}, \emph{error}-type,
	\emph{nl}\}) and \node{Intent} (the change-intent class: fix, feat, refactor, \dots);
	its edges are
	\node{Commit}$\xrightarrow{\text{\edge{MENTIONS}\,(w=\text{TW-IDF})}}$\node{Term} (top
	terms per commit, weighted by TextRank centrality in the commit's word-graph $\times$
	IDF), \node{Term}$\xrightarrow{\text{\edge{COOCCURS}\,(\text{npmi})}}$\node{Term} (the
	global normalised-pointwise-mutual-information term graph),
	\node{Term}$\xrightarrow{\text{\edge{GROUNDS\_IN}}}$\node{ASTNode} (a code term linked
	to the alive identifier leaf it names), and
	\node{Commit}$\xrightarrow{\text{\edge{HAS\_INTENT}}}$\node{Intent}. On Groovy the CSTG
	comprises $20{,}377$ \node{Term} nodes, $188{,}435$ \edge{MENTIONS}, $120{,}340$
	\edge{COOCCURS}, $19{,}315$ \edge{GROUNDS\_IN}, and $8{,}059$ \edge{HAS\_INTENT}.
	Theoretically the layer rests on graph-of-words / TW-IDF (Rousseau \& Vazirgiannis),
	the NPMI corpus term graph (cf.\ TextGCN), defect-semantic typing, and term-risk
	propagation over the NPMI graph. Every CSTG statistic that enters prediction is grown
	online, past-only (see \secref{sec:setup-protocol}); the static term graph and
	term-risk values written to Neo4j are a queryable/visualisation snapshot and are
	\emph{not} read by the online predictor.
	
	\paragraph{Data-quality guarantees.} The finalised graph passes a standing integrity
	audit: no file has more than one structure root, removed nodes always carry a
	\edge{REMOVES} edge, and the online-evolved structure reproduces a fresh
	ground-truth parse of the current file state with precision $=$ recall $=1.0$ on the
	validated sample. A large \emph{ghost-AST} inflation ($\approx\!1.58$\,M nodes left by
	package-restructuring renames outside the labelled set) was identified and corrected,
	so the reported node counts reflect the clean final state. Known, documented
	limitations that a reviewer should be aware of: labels are SZZ-derived and carry the
	usual mining noise; non-\texttt{.java} files are not modelled at the code level; and
	a git rename is currently treated as an add-of-new-path in the delta layer.
	
	\subsection{Inference models}
	\label{sec:setup-models}
	
	The target hardware has \textbf{no GPU}; every method is therefore deliberately
	CPU-only and \textbf{GNN-free}, both to be deployable and to constitute strong,
	honest baselines that any future GNN must beat. The final methodology fixes
	\textbf{five canonical graph-inference paradigms}, chosen because they map one-to-one
	onto the recognised families of graph and knowledge-graph inference, each reads a
	\emph{different, complementary} facet of the structure, and --- crucially --- each is a
	graph \emph{algorithm} that runs unchanged on \emph{any} graph in the family, so all
	five are valid rows across every column of the representation study:
	
	\begin{description}
		\item[RN --- Relational neighbour (wvRN).] Collective classification: a commit's
		score is the label-weighted hub-overlap with past commits (one hop through the shared
		file/developer/term hubs). Complexity $O(\mathrm{nnz})$ of a sparse block matmul.
		\item[PPR --- Personalised PageRank (random walk with restart).] The bipartite
		commit--hub graph is seeded from past buggy vs.\ benign commits and a commit is scored
		by its relative proximity to each class. Complexity $O(\text{iters}\cdot|E|)$ per
		block, run once per class seed set.
		\item[LP --- Label propagation.] Spreading activation of past labels through the
		hubs a few steps, clamped on known past labels. Complexity $O(\text{iters}\cdot|E|)$.
		\item[DW --- DeepWalk-style embedding.] A random-walk / matrix-factorisation node
		embedding (PPMI of the commit--hub co-occurrence, truncated-SVD, the Levy--Goldberg
		form) with a logistic head; marginals from past rows only.
		\item[KGE --- Knowledge-graph embedding (DistMult).] A compact bilinear embedding
		over typed $(\text{commit},\text{rel},\text{hub})$ triples restricted to past
		commits, with commits folded in from their hubs' learned embeddings and a logistic
		head.
	\end{description}
	
	The embedding methods (DW, KGE) are refit on the expanding past window every
	\textbf{$5$ blocks}; RN, PPR, and LP are recomputed each block from the strictly-past
	state. Embedding dimensions are $64$ (DW) and $32$ (KGE). All linear heads are
	balanced logistic regressions.
	
	\paragraph{Two non-graph channels used only in fusion.} Two additional signals are
	\emph{not} admissible as rows of the representation study (they are not graph
	algorithms) but are evaluated as fusion channels:
	\begin{itemize}
		\item \textbf{$M$ --- JIT metrics:} a prequential logistic model over the twelve
		ApacheJIT change metrics; the classical JIT-defect predictor and our metric-only
		reference channel.
		\item \textbf{$G$ --- CSTG feature channel:} a prequential logistic model over the
		CSTG's per-commit features --- the graph-of-words (TextRank-weighted, add/remove
		polarity-tagged) term stream, the NPMI-propagated past-only term-risk prior, a
		$5$-dimensional typed-mass vector over the defect lexicon, and the intent/consistency
		features --- enriched with an online recency / bug-cache block (per-file and
		per-developer time-since-last-buggy touch, recently-buggy touched-file count, re-fix
		flag, commit velocity), each computed from strictly-past state.
	\end{itemize}
	The commit-\emph{message} text channel ($X$) that appeared in earlier iterations is
	\textbf{removed entirely} from the final methodology (it is a strict subset of $G$ and
	adds nothing over it); it is disabled throughout and reported only as a historical
	step.
	
	\paragraph{Fusion.} Method and channel scores are combined by \emph{prequential
		logistic-regression stacking}: at each block a balanced logistic regression is fit on
	the strictly-past predicted probabilities of the constituent models and applied to the
	current block, refit on the expanding window. The \emph{deployed} model is selected by
	an explicit load-vs-performance rule (\secref{sec:setup-fusion-selection}) and is the
	purely graph-native \textbf{$F{+}G = \text{RN}{+}\text{PPR}{+}\text{CSTG}$}, i.e.\ the
	two-method graph fusion $F=\text{RN}{+}\text{PPR}$ augmented with the CSTG semantic
	channel $G$, \emph{without} the non-graph JIT metrics $M$.
	
	\subsection{The two research questions and the graph family}
	\label{sec:setup-rqs}
	
	The experiments answer two questions with one design:
	\textbf{(RQ1, which representation?)} does a richer code/semantic structure improve
	commit-level prediction, and which structural representation is best?
	\textbf{(RQ2, which inference method?)} among typical graph-inference paradigms, which
	best exploits the graph?
	
	Both are answered by a single family of graphs of increasing richness, all sharing the
	Core relational layer, evaluated with the fixed five methods:
	\emph{Core}; \emph{Core$+$AST}; \emph{Core$+$CFG}; \emph{Core$+$DFG};
	\emph{Core$+$PDG}; and the deployed final \emph{Core$+$AST$+$CSTG}. Reading a row
	(a method across graphs) answers RQ1 for that method; reading a column (all methods on
	one graph) answers RQ2 for that graph. A structural layer influences prediction through
	exactly one channel --- the per-commit change-tokens \code{\{edge\}:\{type\}} that
	become typed hubs the walks and embeddings traverse --- so swapping the representation
	means swapping the hub stream while everything else is held fixed. On the final graph,
	the hubs additionally include CSTG \node{Term} nodes via \edge{MENTIONS}, so the same
	five methods traverse the semantic signal too. (Token-seq is included in the
	representation \emph{ablation} but is not part of the final six-graph family, being a
	deliberately weak rival.)
	
	\subsection{Online evaluation protocol}
	\label{sec:setup-protocol}
	
	\paragraph{Prequential, predict-then-grow.} Commits are streamed in chronological
	order (\code{author\_ts}). The protocol is \emph{prequential}: for each commit we
	(i)~predict its label from a model and graph state built only from strictly earlier
	commits, (ii)~reveal the true label, (iii)~grow the graph state (add the commit's
	delta edges, term hubs, recency counters, co-occurrence counts), and (iv)~periodically
	re-learn on the expanded window. This exactly reproduces the information available at
	deployment time and structurally precludes look-ahead leakage.
	
	\paragraph{Warm-up, blocks, and refit schedule.} The first \textbf{$40\%$} of the
	stream ($W = \lfloor 0.40\times 8{,}059\rfloor = 3{,}223$ commits) is a warm-up used
	only to initialise models and is \emph{not} scored; the remaining $4{,}836$ commits are
	the evaluation stream. Predictions are made per commit but models are refreshed on a
	\textbf{block} granularity of $200$ commits (the expanding-window retraining that lets
	the models adapt to concept drift). This warm-up/block schedule, the classifier family,
	and all hyper-parameters are held \emph{identical} across every representation, method,
	and fusion so that the only independent variable in each comparison is the one under
	study.
	
	\paragraph{Leakage controls (stated explicitly).} All scalers, vectorisers, IDF
	weights, NPMI graphs, term-risk priors, embedding bases, and stacking regressions are
	fit on past data only; PPR class seeds never include the commit under test; the
	graph-of-words weighting and recency lookups are commit-local; and the
	\code{HashingVectorizer} used for the CSTG text stream is stateless (a fixed hash with
	no learned vocabulary), so it cannot transport information from a future commit into an
	earlier row. The single-split controlled comparison used against the external text
	baseline (\secref{sec:setup-baselines}) is clearly separated from the deployed online
	path and never mixed with it.
	
	\paragraph{Online operating point.} Because the task is imbalanced and drifting, the
	hard-decision threshold is itself tuned \emph{online}: it is re-tuned on the past-only
	predictions every $150$ commits to maximise buggy-F1 on already-seen data, giving a
	leakage-free deployment operating point at which all threshold-based metrics are read.
	Rolling trajectories are computed over a trailing window of $\text{ROLL}=800$ commits
	(a visualisation smoothing only --- it never enters the model or the reported scalars,
	which are computed once over the entire evaluation stream).
	
	\subsection{Evaluation metrics}
	\label{sec:setup-metrics}
	
	Because the positive class is rare and drifting, no single metric suffices; we report a
	\textbf{seven-metric} suite plus two effort-aware measures, and treat the
	threshold-free and class-balanced metrics as primary.
	
	\begin{itemize}
		\item \textbf{Threshold-free ranking:} \emph{ROC-AUC} (a.k.a.\ AUC) and, where
		reported, \emph{PR-AUC} (area under precision--recall; its random baseline equals the
		positive rate, which is why absolute PR-AUC falls in the low-defect tail even for a
		good model).
		\item \textbf{Operating-point metrics at the online-tuned threshold:} \emph{Precision},
		\emph{Recall} (buggy-class sensitivity), \emph{Buggy-F1} (F1 of the positive class),
		\emph{Macro-F1} (unweighted mean of both classes' F1 --- robust to imbalance),
		\emph{G-Mean} ($\sqrt{\text{TPR}\cdot\text{TNR}}$), and \emph{MCC} (Matthews
		correlation --- the single most reliable summary under imbalance).
		\item \textbf{Effort-aware (deployment triage) metrics:} normalised
		$P_{\mathrm{opt}}$ (Kamei et al.: how close the model's inspection order is to the
		ideal buggy-first, least-effort-first order, normalised so $0.5$ is random and $1.0$
		optimal) and \emph{ACC@20\%LOC} (recall achieved when only the top $20\%$ of churned
		lines are inspected). Effort is measured as \code{la}$+$\code{ld}.
	\end{itemize}
	The two \emph{class-balanced} metrics --- Macro-F1 and G-Mean --- are the ones we
	prioritise for model selection, precisely because they do not let a model win by
	exploiting the majority class in the low-defect regime. We emphasise that, given the
	concept drift, the meaningful reading of an online trajectory is \emph{relative} (which
	method/graph stays highest) rather than the absolute end-of-stream value, since the
	falling positive rate depresses every method's tail together.
	
	\subsection{Fusion selection rule}
	\label{sec:setup-fusion-selection}
	
	The deployed fusion is not hand-picked. We evaluate \emph{all} $2^5-1=31$ combinations
	of the five methods on the final graph, score each on the full seven-metric suite, and
	select the smallest combination whose balanced score
	$(\text{Macro-F1}+\text{G-Mean})/2$ is within a fixed tolerance ($0.005$) of the best ---
	a transparent trade-off of \emph{load} (fewer methods, cheaper online) against
	\emph{performance}. This selects $F=\text{RN}{+}\text{PPR}$. We then test, on top of
	$F$, whether the explicit CSTG channel $G$ and/or the non-graph JIT metrics $M$ add
	anything: $G$ helps markedly and $M$ does not (adding $M$ alone hurts), so the deployed
	model is $F{+}G=\text{RN}{+}\text{PPR}{+}\text{CSTG}$. Reporting the full combinatorial
	ablation, and choosing by an explicit rule, is what makes the deployed model
	defensible rather than cherry-picked.
	
	\subsection{Baselines and comparison protocols}
	\label{sec:setup-baselines}
	
	We compare against baselines at two levels.
	\textbf{(i) Within-project JIT baselines:} logistic regression, random forest, and
	gradient boosting over the twelve ApacheJIT metrics --- the standard JIT-defect
	predictor --- evaluated under the identical online protocol, plus a naive prior-rate
	predictor. \textbf{(ii) A strong external text baseline:} a $768$-dimensional
	diff-text TF-IDF model with a random forest, the strongest published lexical baseline.
	Because that baseline was reported on a single chronological split (train first $70\%$,
	test last $15\%$; a hard, low-defect tail with $\approx\!7.5\%$ positives), we
	reproduce that \emph{exact} split and test set for a like-for-like comparison, using a
	matched classifier so any difference comes from the features and not the learner; these
	single-split numbers are reported separately from, and never conflated with, the
	fully-online results. As a sanity check our JIT-metrics model reproduces the baseline
	logistic regression on that split to three decimals, confirming the protocols are
	aligned. For the CSTG internal ablation specifically, we additionally report a
	\emph{same-learner} (matched random forest) plain-TF-IDF row so the ``does the CSTG
	representation beat a bag-of-words'' comparison is under one learner rather than mixing
	linear and non-linear models.
	
	\subsection{Scalability, complexity, and statistical-significance setup}
	\label{sec:setup-scalability}
	
	Beyond accuracy, the paper characterises the system's cost and the statistical
	reliability of its comparisons. All of this is computed \emph{without rebuilding} the
	knowledge graph: final sizes and growth come from read-only queries against the built
	graph and from attributing delta edges to their source commits in online order;
	prediction latency is measured by replaying the cached graph family in memory; and
	per-commit build/modify cost is obtained by an instrumented replay of the real
	parse-and-diff work on a stratified sample of commits, cross-checked against an
	analytical operation-count proxy.
	
	\begin{itemize}
		\item \textbf{Statistical profile.} Per-layer node/edge counts, type and token
		vocabularies, per-file size distributions (with Gini and skew), and a discrete
		power-law fit of the per-commit change-size distribution (heavy-tailed, exponent
		$2<\alpha<3$ on every layer).
		\item \textbf{Per-commit growth and modification.} Cumulative nodes/edges over the
		commit stream and over calendar time (overlaid on the drifting bug rate), the
		per-commit change-size distribution, and the buggy-vs-benign change-size separation
		--- which is large and highly significant (Mann--Whitney $p\ll 10^{-100}$ on every
		layer; buggy commits make on the order of $2$--$3\times$ larger structural changes),
		the basic reason the graph is predictive.
		\item \textbf{Time and space complexity.} Build/modify cost is shown to be near-linear
		in file size ($R^2\approx0.99$), i.e.\ per-commit work is $O(\text{change})$ and space
		is $O(\text{alive nodes}+\text{retained history})$; prediction latency scales with
		graph density, and the deployed graph-native fusion predicts at sub-millisecond per
		commit (thousands of commits/second), so the whole pipeline is comfortably real-time
		and GPU-free. Each per-commit method's cost is annotated with its known asymptotic
		complexity and validated against the measured block-time-vs-graph-size growth.
		\item \textbf{Statistical significance.} Headline comparisons are tested on paired
		per-commit out-of-sample scores: DeLong's test for ROC-AUC differences (Holm-corrected
		within each family of comparisons), paired bootstrap tests and $95\%$ confidence
		intervals for PR-AUC/MCC/G-Mean, and Cliff's~$\delta$ effect sizes. This lets us state
		which accuracy gaps are significant (e.g.\ graph richness significantly helps the
		informative methods; PPR is significantly the best method on the final graph; the
		semantic channel $F{+}G$ significantly improves every headline metric over $F$) and
		which are within noise.
	\end{itemize}
	
	\subsection{Implementation and reproducibility}
	\label{sec:setup-impl}
	
	The graph store is Neo4j (final on-disk footprint $\approx\!3.3$\,GB; working footprint
	$\approx\!5.6$\,GB with transaction logs). Java is parsed with a single deterministic
	parser (\texttt{javalang}) for the whole pipeline, in killable subprocesses with a
	watchdog so a pathological file cannot stall the build; a persistent parser worker and
	an in-memory before-graph cache keep the online build tractable, and a range index on
	each node label's \code{id} keeps per-commit cost constant as the graph grows (an
	absent index was, at one point, the dominant cost --- turning an $O(\text{change})$
	build into an accidentally quadratic one). All inference code is CPU-only Python
	(scikit-learn linear models, custom sparse graph routines, gensim/PyKEEN for the
	embeddings where used); no GNN and no GPU are required. Feature streams and per-method
	scores are cached to disk so every experiment, table, figure, and statistical test can
	be regenerated without a live database or a graph rebuild, and companion Jupyter
	notebooks reproduce all figures from those caches. Randomised components (embedding
	initialisation, bootstrap resampling) use fixed seeds. Determinism of node identity
	across versions (the same source always yields the same parse-local ids) is what makes
	the online build itself reproducible.
	
	\paragraph{Summary of the final configuration.} The deployed predictor is a
	fully-online, GPU-free, multimodal commit knowledge graph on \texttt{apache/groovy}
	($8{,}059$ ApacheJIT commits, $\approx\!20\%$ buggy, 2003--2019): a Core process layer,
	an AST delta layer, and the CSTG semantic-text layer, over which the fusion
	$F{+}G=\text{RN}{+}\text{PPR}{+}\text{CSTG}$ predicts each commit prequentially from
	strictly-past state, warm-up $40\%$ ($3{,}223$ commits), block $200$, evaluated on the
	remaining $4{,}836$ commits with a seven-metric plus effort-aware suite and paired
	significance testing.

	% =====================================================================
	\section{Environment, Software Versions, and Model-Configuration Checklist}
	\label{sec:setup-checklist}
	% =====================================================================
	For completeness and reproducibility, this section answers a standard checklist of
	setup details --- operating system and hardware, software and library versions,
	training hyper-parameters, model architecture, and stopping/tuning. Several checklist
	items are framed for \emph{deep-learning} pipelines (optimiser, learning rate, batch
	size, epochs, hidden-layer dimensions, early stopping). Our final methodology is
	deliberately \textbf{GPU-free and GNN-free}: the inference models are classical
	graph algorithms (random-walk / label-propagation / matrix-factorisation) with
	\emph{linear} (logistic-regression) heads, plus one compact bilinear knowledge-graph
	embedding trained by a short custom SGD. Where a checklist item has no analogue in
	this design, we say so explicitly and explain why, rather than inventing a value.

	\subsection{Operating system and hardware}
	\label{sec:setup-hw}
	All building and all experiments were run on a single commodity laptop, with no GPU
	and no cluster:
	\begin{itemize}[nosep,leftmargin=1.4em]
		\item \textbf{Operating system:} Microsoft Windows~11 (64-bit, build $26{,}200$).
		The pipeline is OS-agnostic (pure Python plus a Dockerised Neo4j) and also runs on
		Linux/macOS; the Neo4j administration helpers assume a Docker-hosted server.
		\item \textbf{CPU:} Intel Core i7-13620H (13th generation), $10$ physical cores /
		$16$ logical threads, $2.4$\,GHz base clock. \textbf{No GPU is used at any stage}
		(build, inference, or analysis).
		\item \textbf{Memory (RAM):} $16$\,GB. The whole pipeline is designed to run within
		this budget: the online build stores only one evolving structure per file (not a
		snapshot per commit), and the inference/analysis code streams and caches to disk, so
		no stage requires loading the full graph into memory at once.
		\item \textbf{Graph store:} Neo4j runs in a Docker container; the built Groovy graph
		occupies $\approx\!3.3$\,GB on disk ($\approx\!5.6$\,GB with transaction logs).
	\end{itemize}

	\subsection{Software and library versions}
	\label{sec:setup-versions}
	The exact versions used are listed in Table~\ref{tab:setup-versions}. The final
	pipeline depends only on the mainstream scientific-Python stack, a Java parser, and a
	Neo4j driver; there is \textbf{no deep-learning framework} (no PyTorch, TensorFlow, or
	Keras) anywhere on the critical path, consistent with the CPU-only design.

	\begin{table}[t]
		\centering
		\caption{Software and library versions used for all builds and experiments.}
		\label{tab:setup-versions}
		\begin{tabular}{lll}
			\toprule
			Component & Version & Role \\
			\midrule
			Python            & 3.13.5   & runtime for the entire pipeline \\
			Neo4j             & 2026.05 (Community) & graph store (Dockerised) \\
			\code{neo4j} (driver) & 6.2.0 & Python$\to$Neo4j Bolt driver \\
			git               & 2.45.1   & reads per-commit file blobs / diffs \\
			javalang          & 0.13.0   & Java parser (AST/CFG/DFG/PDG/SEQ builders) \\
			NumPy             & 2.2.6    & arrays; the graph-method kernels \\
			SciPy             & 1.16.2   & sparse matrices; statistics (Mann--Whitney, ranks) \\
			scikit-learn      & 1.7.2    & logistic-regression heads, TF-IDF, RF/GB baselines \\
			pandas            & 2.3.2    & CSV / tabular handling \\
			networkx          & 3.5      & in-memory graph construction in the builders \\
			matplotlib        & 3.10.6   & figures \\
			nbformat / nbconvert & 5.10.4 / 7.17.1 & build and execute the result notebooks \\
			PyYAML            & 6.0.2    & per-project configuration (reproducibility package) \\
			\bottomrule
		\end{tabular}
	\end{table}

	\paragraph{A note on gensim / PyKEEN.} Earlier iterations experimented with
	\code{gensim} (node2vec/DeepWalk) and \code{PyKEEN} (TransE/DistMult). These are
	\textbf{not} used by the final methodology: the deployed DeepWalk embedding (DW) is a
	self-contained truncated-SVD of a PPMI matrix, and the deployed knowledge-graph
	embedding (KGE) is our own compact DistMult in NumPy (below). No external embedding
	library is on the final critical path.

	\subsection{Training hyper-parameters}
	\label{sec:setup-hparams}
	The checklist items \emph{optimiser}, \emph{learning rate}, \emph{batch size}, and
	\emph{number of epochs} are properties of gradient-descent neural training. Our final
	models are not neural, so most of these \textbf{do not apply}; the few genuine
	training knobs are given precisely.

	\begin{description}[leftmargin=1.4em]
		\item[Optimiser.] \emph{Mostly not applicable.} The linear heads on every graph
		method and on the $M$ and $G$ fusion channels are \textbf{balanced logistic
		regressions solved by L-BFGS} (scikit-learn, \code{solver="lbfgs"},
		\code{class\_weight="balanced"}, \code{max\_iter}$=1000$--$2000$); L-BFGS is a
		quasi-Newton batch solver, so there is no stochastic optimiser, no learning-rate
		schedule, and no mini-batching. The \emph{only} stochastic-gradient component in the
		whole system is the DistMult KGE, trained with plain SGD.
		\item[Learning rate.] \emph{Only the KGE has one:} the DistMult embedding uses a
		fixed SGD learning rate of $\eta=0.05$. All logistic-regression heads are L-BFGS and
		have no learning rate. DeepWalk (DW) is closed-form (SVD) and has none.
		\item[Batch size.] \emph{Not applicable in the neural sense.} There is no
		mini-batch training. The one gradient method (KGE) processes its past triples in
		internal chunks of $2{,}048$ for vectorised updates, but this is an implementation
		detail of a full-pass SGD, not a tuned batch-size hyper-parameter. The evaluation
		protocol does have a \emph{block} granularity ($200$ commits) at which models are
		refreshed, but that is a streaming-protocol constant (\secref{sec:setup-protocol}),
		not a training batch size.
		\item[Number of epochs.] \emph{Only the KGE has epochs:} the DistMult SGD runs
		$3$ epochs over the strictly-past triples at each refit; negative sampling uses
		$3$ corrupted tails per positive. No other model has an epoch count: the
		logistic heads run to L-BFGS convergence (capped at \code{max\_iter}), and DW/RN/PPR/LP
		are single-pass or fixed-iteration (below), not epoch-based.
	\end{description}

	\noindent The fixed-iteration counts of the non-embedding graph methods (which play the
	role ``epochs'' would in a neural model) are: \textbf{PPR} --- $40$ power-iterations of
	random-walk-with-restart, restart probability $\alpha=0.15$, run once per class seed set
	per block; \textbf{LP} --- $3$ propagation steps with clamping on known past labels;
	\textbf{RN} --- a single one-hop weighted vote (no iteration). These are held fixed
	across all projects and all graphs.

	\subsection{Model architecture}
	\label{sec:setup-arch}
	The checklist item \emph{hidden-layer dimensions} presupposes a multi-layer neural
	network. Our models have \textbf{no hidden layers} --- there is no multi-layer
	network, and specifically \textbf{no GNN} (a heterogeneous GNN is explicitly scoped as
	future work, not part of this study). The only learned latent dimensionalities are the
	two shallow embedding sizes already reported: \textbf{DeepWalk (DW) $=64$} and
	\textbf{DistMult KGE $=32$}. Beyond those:
	\begin{itemize}[nosep,leftmargin=1.4em]
		\item \textbf{RN, PPR, LP} produce a scalar risk score directly from the
		commit--hub incidence / transition matrix; they have no latent layer at all.
		\item \textbf{DW} embeds commits into $\mathbb{R}^{64}$ by truncated-SVD of the PPMI
		of the commit--hub co-occurrence (Levy--Goldberg form), then a logistic head; there
		is no hidden layer, only this one factorisation dimension.
		\item \textbf{KGE (DistMult)} embeds hubs and relations into $\mathbb{R}^{32}$ with a
		bilinear score; commits are folded in as the mean of their hubs' embeddings, then a
		logistic head. Again a single latent dimension, no hidden layers.
		\item \textbf{Text/structural feature widths} (not ``hidden'' layers, but the input
		dimensionalities worth recording): the structural TF--IDF and the CSTG text stream
		are hashed to $2^{18}$ features (a fixed, stateless hashing vectoriser); the CSTG
		typed-mass vector is $5$-dimensional; the JIT-metric channel is $12$-dimensional.
	\end{itemize}
	So the ``architecture'' of the deployed model is: five shallow graph-inference scorers
	(two of them with a $64$- or $32$-dim embedding) whose scalar outputs are combined by a
	single balanced logistic-regression stacker --- no hidden layers, no GNN.

	\subsection{Stopping and tuning}
	\label{sec:setup-stopping}
	\begin{description}[leftmargin=1.4em]
		\item[Early-stopping criteria.] \emph{Not used, and not applicable.} Because the
		models are non-neural, there is no validation-loss-based early stopping. The KGE
		SGD runs a fixed $3$ epochs; the logistic heads stop at L-BFGS convergence (or the
		\code{max\_iter} cap); the fixed-iteration methods (PPR $40$, LP $3$) run to their
		set iteration count. The streaming protocol itself has no ``stopping'': every
		post-warm-up commit is scored, and the run ends only when the commit stream ends.
		\item[Hyper-parameter tuning strategy.] We deliberately \textbf{do not tune the
		model hyper-parameters per project or per graph}, and this is a design choice, not
		an omission: the whole study is a \emph{controlled comparison}, so the protocol
		constants (warm-up $40\%$, block $200$, roll $800$, refit every $5$ blocks), the
		embedding sizes ($64$/$32$), the iteration counts (PPR $40$, LP $3$, KGE
		$3$ epochs / $\eta=0.05$ / $3$ negatives), and the classifier family (balanced L-BFGS
		logistic regression) are \textbf{held fixed across every representation, method,
		fusion, and project}. Fixing them is what makes cross-configuration and
		cross-project differences attributable to the factor under study rather than to
		per-cell tuning. The one thing that \emph{is} selected --- the deployed fusion
		formula --- is chosen not by a hyper-parameter search but by the transparent,
		reported combinatorial rule of \secref{sec:setup-fusion-selection} (evaluate all
		$2^5-1$ method combinations, take the smallest within a fixed tolerance of the best
		balanced score). The threshold at which hard decisions are read is likewise not a
		tuned constant but is re-derived \emph{online} from past-only data every $150$
		commits (\secref{sec:setup-protocol}). The default hyper-parameter values were set
		once, up front, from standard choices in the literature (e.g.\ $\alpha=0.15$ for
		PPR) and small sanity checks; there was no automated grid/random/Bayesian search
		over them, precisely to avoid over-fitting the evaluation.
	\end{description}

	\paragraph{Cross-project note.} Every constant above is identical across all subject
	projects (Groovy and the additional Apache projects processed with the reproducibility
	package); only the data changes. The single genuinely per-project quantity is the base
	commit that seeds the structural layers (\secref{sec:setup-data}), which is a property
	of each repository's history, not a tuned hyper-parameter.

\end{document} from the parent .tex.
% =====================================================================
\documentclass[11pt,a4paper]{article}
\usepackage[margin=2.4cm]{geometry}
\usepackage[T1]{fontenc}
\usepackage[utf8]{inputenc}
\usepackage{lmodern}
\usepackage{amsmath,amssymb}
\usepackage{booktabs}
\usepackage{array}
\usepackage{enumitem}
\usepackage{xcolor}
\usepackage[hidelinks]{hyperref}

% ---- macros used throughout (match the main document's conventions) ----
\providecommand{\code}[1]{\texttt{#1}}
\providecommand{\node}[1]{\textsf{#1}}
\providecommand{\edge}[1]{\textsc{#1}}

% ---- helper for safe section references (avoids any parsing issues with \S\ref) ----
\newcommand{\secref}[1]{\S\ref{#1}}

\begin{document}
	
	\section{Experimental Setup}
	\label{sec:experimental-setup}
	
	This section specifies, in full, the data, the knowledge-graph construction, the
	inference models, the online evaluation protocol, the metrics, the baselines, and
	the reproducibility conditions under which every result in this paper was produced.
	Two principles govern the whole design and are worth stating up front because they
	recur throughout. \textbf{First, everything is \emph{just-in-time (JIT) and fully
			online}:} the graph grows one commit at a time in true chronological order, and a
	commit is always predicted from a model and a graph state built \emph{only} from
	strictly earlier commits --- there is no random split, no look-ahead, and no
	statistic is ever fit on data that includes or postdates the commit under test.
	\textbf{Second, the study is designed as a controlled comparison:} when we ask which
	code representation or which inference method is best, exactly one factor is varied
	at a time while the protocol, the classifier family, the warm-up, and the block
	schedule are held fixed, so any measured difference is attributable to that factor
	and not to a confound.
	
	\subsection{Subject program, commit scope, and labels}
	\label{sec:setup-data}
	
	\paragraph{Project.} All experiments are conducted on \texttt{apache/groovy}, a
	mature, long-lived Java project (the Groovy programming language and its compiler,
	runtime, and tooling). We work from a full local \texttt{git} clone, so the raw
	material for every commit --- its message, its parent(s), and the exact byte content
	of every file before and after the change --- is available offline and deterministic.
	
	\paragraph{Labelled commit scope.} The study set is the \textbf{$8{,}059$ commits}
	of \texttt{apache/groovy} that appear in the \textbf{ApacheJIT} dataset
	(\texttt{data/apachejit/projects/apache\_groovy.csv}), spanning \textbf{2003--2019}.
	ApacheJIT provides, for each commit, a \emph{buggy/benign} label derived by an
	SZZ-style procedure (a commit is buggy if a later fix commit points back to it) and
	the twelve standard JIT change metrics. These $8{,}059$ commits are the only ones we
	treat as the labelled study set and the only ones grown into the AST/structural and
	semantic layers. The surrounding history is retained in the database as
	\node{Commit} nodes for context (the full clone has $16{,}887$ commits), but only the
	ApacheJIT commits are flagged \code{in\_jit=true} and are the units of prediction.
	
	\paragraph{Class balance and its consequences.} The labelled set is imbalanced:
	\textbf{$1{,}614$ buggy / $6{,}445$ benign}, i.e.\ a $\approx\!20\%$ positive rate
	overall. Crucially, the positive rate is \emph{not stationary}: it rises from
	$\approx\!6\%$ in 2003 to a peak of $\approx\!39\%$ in 2012 and then falls back to
	$\approx\!6\%$ by 2019 (measured directly from the labels). This pronounced
	\emph{concept drift} is the single most important property of the data for our
	purposes: it is why a random or single-split evaluation is inappropriate, why we
	adopt a time-aware online protocol, and why the tail of every online trajectory
	declines (the achievable precision/recall shrink as positives become rare). We
	return to this when defining the protocol (\secref{sec:setup-protocol}) and the
	metrics (\secref{sec:setup-metrics}).
	
	\paragraph{Base snapshot.} The structural layers are seeded once at a fixed
	\emph{base} commit --- \texttt{408b29851d7bbe4d343340832297e4be7e0c5578}, the first
	substantial Java commit --- whose $120$ \texttt{.java} files (${\approx}46$k AST
	nodes) constitute the initial code structure onto which all subsequent per-commit
	changes are laid. The base is a property of the \emph{code}, not of the labelled
	set; it merely fixes the starting state of the evolving graph.
	
	\subsection{Knowledge-graph construction and online growth}
	\label{sec:setup-kg}
	
	The knowledge graph is a \emph{heterogeneous information network} (HIN): typed nodes
	and typed edges in one store (Neo4j). It fuses three families of evidence that the
	defect-prediction literature usually treats separately --- \emph{process} evidence
	(who changed what, when, and how often), \emph{code-structure} evidence (the
	syntactic/semantic shape of the change and its surroundings), and \emph{textual}
	evidence (the natural-language and lexical content of the change). The design
	rationale for a graph (rather than a flat feature table) is that anchoring each
	commit's change to \emph{shared, persistent} code entities is what makes relational
	and neighbourhood reasoning across the entire history possible.
	
	\paragraph{Core (process / relational) layer.} The Core layer is the classical
	commit graph and is present in \emph{every} experimental configuration. Its node
	types are \node{Project}, \node{Branch}, \node{Developer}, \node{Commit},
	\node{File}, and \node{Issue}; its edges encode authorship
	(\node{Commit}$\xrightarrow{\text{\edge{AUTHORED\_BY}}}$\node{Developer}), the
	history spine (\node{Commit}$\xrightarrow{\text{\edge{PARENT\_OF}}}$\node{Commit}),
	file changes
	(\edge{ADDED}/\edge{MODIFIED}/\edge{DELETED}/\edge{RENAMED\_FROM}/\edge{RENAMED\_TO}
	between \node{Commit} and \node{File}), and issue links
	(\node{Commit}$\xrightarrow{\text{\edge{FIXES\_ISSUE}}}$\node{Issue}). Every
	\node{Commit} also carries the twelve ApacheJIT change metrics (size \code{la},
	\code{ld}; dispersion \code{nf}, \code{nd}, \code{ns}, \code{ent}; history and
	experience \code{age}, \code{ndev}, \code{nuc}, \code{aexp}, \code{arexp},
	\code{asexp}) and its \code{author\_ts} (the UNIX timestamp that defines online
	arrival order), \code{buggy} label, and \code{jit\_year}.
	
	\paragraph{Structural (code) layers as \emph{delta} graphs.} On top of the Core, we
	materialise the code structure of every changed \texttt{.java} file. The key modelling
	decision --- and a central contribution --- is that a commit is not attached to a
	per-commit snapshot blob; it is attached, by four typed \emph{delta} edges
	(\edge{ADDS}, \edge{REMOVES}, \edge{UPDATES}, \edge{MOVES}), directly to the specific
	code nodes it changed, laid on top of a single evolving structure per file. Concretely,
	for each tracked file the graph holds one \emph{current} structure; when a commit
	touches the file we (i) parse the before- and after-states with the \emph{same}
	deterministic builder, (ii) diff them, (iii) write the delta edges from the
	\node{Commit} to the affected nodes, and (iv) evolve the stored structure so the next
	commit diffs against the correct state. Node identity is threaded across versions by a
	deterministic, per-file parse-local-id map, so the ``after'' map of commit $N$ is
	byte-identical to the ``before'' map of the next commit that touches the file when the
	intervening source is unchanged. Removed nodes are \emph{retained} and flagged
	(\code{alive=false}, \code{removed\_by}), so the complete change history is queryable
	while the live structure always reflects the current code. This makes storage grow
	with the \emph{amount of change}, not with history-length $\times$ file-size, and it is
	what makes the whole graph online-extensible.
	
	We instantiate this delta-graph scheme for \textbf{five} structural representations of
	increasing coarseness, all built per changed \texttt{.java} file and attached at the
	\node{File} level:
	\begin{itemize}
		\item \textbf{AST} --- the abstract syntax tree; one node per syntactic construct
		(declaration, statement, expression, operator, literal, and each identifier/operator
		leaf). Finest granularity ($136$ node types).
		\item \textbf{CFG} --- the intra-procedural control-flow graph; one node per statement
		or control construct ($18$ node types).
		\item \textbf{DFG} --- the def--use data-flow graph over local variables; one node per
		definition/use site ($7$ node types).
		\item \textbf{PDG} --- the program-dependence graph; the CFG's statement vertices plus
		control-dependence and data-dependence \emph{edges} (no new nodes; $18$ node types).
		We label it PDG rather than CPG because we deliberately keep dependence at statement
		granularity and do \emph{not} embed the AST subtree under each statement; a canonical
		code property graph would contain the AST as a subgraph and is out of scope.
		\item \textbf{Token-seq} --- a deliberately weak flat statement sequence (statement
		nodes chained by \edge{NEXT}, no control or data structure); a lower-bound rival used
		only in the representation comparison.
	\end{itemize}
	Table~\ref{tab:setup-layers} reports the materialised scale of each layer. The
	${\approx}10\times$ gap between AST ($3.70$\,M nodes) and the others ($0.26$--$0.47$\,M)
	is a pure \emph{granularity} effect: a statement such as \code{int x = foo(a,b)+1;} is
	one node in CFG/PDG but a dozen-plus nodes in the AST. CFG and PDG have identical node
	counts because the PDG adds only dependence edges to the CFG's vertex set.
	
	\begin{table}[t]
		\centering
		\caption{Materialised scale of each structural layer and its per-commit change-token
			stream (Groovy, $8{,}059$ commits). ``Token types'' counts the distinct
			\code{\{edge\}:\{node-type\}} change-tokens; ``Commits w/ tok.'' is the number of
			labelled commits that emit at least one token in that layer.}
		\label{tab:setup-layers}
		\begin{tabular}{lrrrrr}
			\toprule
			Layer & Nodes & Delta-edges & Node types & Token types & Commits w/ tok. \\
			\midrule
			AST       & 3{,}699{,}439 & 3{,}614{,}922 & 136 & 314 & 7{,}523 \\
			CFG       &   372{,}488 &   408{,}301 & 18 & 53 & 6{,}194 \\
			DFG       &   469{,}662 &   648{,}553 &  7 & 27 & 6{,}066 \\
			PDG       &   372{,}488 &   411{,}869 & 18 & 63 & 6{,}261 \\
			Token-seq &   257{,}708 &   293{,}287 & 16 & 47 & 6{,}011 \\
			\bottomrule
		\end{tabular}
	\end{table}
	
	\paragraph{The differ, and why matching is fair across layers.} The AST layer is
	diffed as a \emph{tree} by bottom-up subtree hashing (identical subtrees in the before
	and after trees are matched greedily, preferring equal source positions; unmatched
	before-nodes are deletions, unmatched after-nodes insertions, matched-but-relabelled
	nodes updates, reparented nodes moves). The four non-tree layers (CFG/DFG/PDG/Token-seq)
	are general typed graphs and are diffed by a single \emph{node-identity-lite} matcher
	applied \emph{identically} to each: matching is scoped within a method (\code{class::name},
	stable across commits) and keyed by \emph{ordinal} (the $k$-th node of a given type in a
	method maps to the $k$-th such node in the other version), which is robust to line
	shifts. Using one shared differ for the alternatives is the fairness guarantee for the
	representation comparison: only the subgraph varies, never the difference model. We
	disclose that the AST uses a tree differ while the others use the ordinal matcher --- a
	deliberate specialisation (trees vs.\ general graphs), not a hidden confound.
	
	\paragraph{Semantic-text layer (CSTG).} The third first-class layer is the
	\textbf{Commit Semantic-Text Graph (CSTG)}, which models the \emph{textual}
	description of a change (commit message $+$ diff text) as a graph grounded in code.
	Its node types are \node{Term} (a typed vocabulary term, with \code{kind} $\in$
	\{\emph{code}-identifier, \emph{bug}-lexicon, \emph{action}, \emph{error}-type,
	\emph{nl}\}) and \node{Intent} (the change-intent class: fix, feat, refactor, \dots);
	its edges are
	\node{Commit}$\xrightarrow{\text{\edge{MENTIONS}\,(w=\text{TW-IDF})}}$\node{Term} (top
	terms per commit, weighted by TextRank centrality in the commit's word-graph $\times$
	IDF), \node{Term}$\xrightarrow{\text{\edge{COOCCURS}\,(\text{npmi})}}$\node{Term} (the
	global normalised-pointwise-mutual-information term graph),
	\node{Term}$\xrightarrow{\text{\edge{GROUNDS\_IN}}}$\node{ASTNode} (a code term linked
	to the alive identifier leaf it names), and
	\node{Commit}$\xrightarrow{\text{\edge{HAS\_INTENT}}}$\node{Intent}. On Groovy the CSTG
	comprises $20{,}377$ \node{Term} nodes, $188{,}435$ \edge{MENTIONS}, $120{,}340$
	\edge{COOCCURS}, $19{,}315$ \edge{GROUNDS\_IN}, and $8{,}059$ \edge{HAS\_INTENT}.
	Theoretically the layer rests on graph-of-words / TW-IDF (Rousseau \& Vazirgiannis),
	the NPMI corpus term graph (cf.\ TextGCN), defect-semantic typing, and term-risk
	propagation over the NPMI graph. Every CSTG statistic that enters prediction is grown
	online, past-only (see \secref{sec:setup-protocol}); the static term graph and
	term-risk values written to Neo4j are a queryable/visualisation snapshot and are
	\emph{not} read by the online predictor.
	
	\paragraph{Data-quality guarantees.} The finalised graph passes a standing integrity
	audit: no file has more than one structure root, removed nodes always carry a
	\edge{REMOVES} edge, and the online-evolved structure reproduces a fresh
	ground-truth parse of the current file state with precision $=$ recall $=1.0$ on the
	validated sample. A large \emph{ghost-AST} inflation ($\approx\!1.58$\,M nodes left by
	package-restructuring renames outside the labelled set) was identified and corrected,
	so the reported node counts reflect the clean final state. Known, documented
	limitations that a reviewer should be aware of: labels are SZZ-derived and carry the
	usual mining noise; non-\texttt{.java} files are not modelled at the code level; and
	a git rename is currently treated as an add-of-new-path in the delta layer.
	
	\subsection{Inference models}
	\label{sec:setup-models}
	
	The target hardware has \textbf{no GPU}; every method is therefore deliberately
	CPU-only and \textbf{GNN-free}, both to be deployable and to constitute strong,
	honest baselines that any future GNN must beat. The final methodology fixes
	\textbf{five canonical graph-inference paradigms}, chosen because they map one-to-one
	onto the recognised families of graph and knowledge-graph inference, each reads a
	\emph{different, complementary} facet of the structure, and --- crucially --- each is a
	graph \emph{algorithm} that runs unchanged on \emph{any} graph in the family, so all
	five are valid rows across every column of the representation study:
	
	\begin{description}
		\item[RN --- Relational neighbour (wvRN).] Collective classification: a commit's
		score is the label-weighted hub-overlap with past commits (one hop through the shared
		file/developer/term hubs). Complexity $O(\mathrm{nnz})$ of a sparse block matmul.
		\item[PPR --- Personalised PageRank (random walk with restart).] The bipartite
		commit--hub graph is seeded from past buggy vs.\ benign commits and a commit is scored
		by its relative proximity to each class. Complexity $O(\text{iters}\cdot|E|)$ per
		block, run once per class seed set.
		\item[LP --- Label propagation.] Spreading activation of past labels through the
		hubs a few steps, clamped on known past labels. Complexity $O(\text{iters}\cdot|E|)$.
		\item[DW --- DeepWalk-style embedding.] A random-walk / matrix-factorisation node
		embedding (PPMI of the commit--hub co-occurrence, truncated-SVD, the Levy--Goldberg
		form) with a logistic head; marginals from past rows only.
		\item[KGE --- Knowledge-graph embedding (DistMult).] A compact bilinear embedding
		over typed $(\text{commit},\text{rel},\text{hub})$ triples restricted to past
		commits, with commits folded in from their hubs' learned embeddings and a logistic
		head.
	\end{description}
	
	The embedding methods (DW, KGE) are refit on the expanding past window every
	\textbf{$5$ blocks}; RN, PPR, and LP are recomputed each block from the strictly-past
	state. Embedding dimensions are $64$ (DW) and $32$ (KGE). All linear heads are
	balanced logistic regressions.
	
	\paragraph{Two non-graph channels used only in fusion.} Two additional signals are
	\emph{not} admissible as rows of the representation study (they are not graph
	algorithms) but are evaluated as fusion channels:
	\begin{itemize}
		\item \textbf{$M$ --- JIT metrics:} a prequential logistic model over the twelve
		ApacheJIT change metrics; the classical JIT-defect predictor and our metric-only
		reference channel.
		\item \textbf{$G$ --- CSTG feature channel:} a prequential logistic model over the
		CSTG's per-commit features --- the graph-of-words (TextRank-weighted, add/remove
		polarity-tagged) term stream, the NPMI-propagated past-only term-risk prior, a
		$5$-dimensional typed-mass vector over the defect lexicon, and the intent/consistency
		features --- enriched with an online recency / bug-cache block (per-file and
		per-developer time-since-last-buggy touch, recently-buggy touched-file count, re-fix
		flag, commit velocity), each computed from strictly-past state.
	\end{itemize}
	The commit-\emph{message} text channel ($X$) that appeared in earlier iterations is
	\textbf{removed entirely} from the final methodology (it is a strict subset of $G$ and
	adds nothing over it); it is disabled throughout and reported only as a historical
	step.
	
	\paragraph{Fusion.} Method and channel scores are combined by \emph{prequential
		logistic-regression stacking}: at each block a balanced logistic regression is fit on
	the strictly-past predicted probabilities of the constituent models and applied to the
	current block, refit on the expanding window. The \emph{deployed} model is selected by
	an explicit load-vs-performance rule (\secref{sec:setup-fusion-selection}) and is the
	purely graph-native \textbf{$F{+}G = \text{RN}{+}\text{PPR}{+}\text{CSTG}$}, i.e.\ the
	two-method graph fusion $F=\text{RN}{+}\text{PPR}$ augmented with the CSTG semantic
	channel $G$, \emph{without} the non-graph JIT metrics $M$.
	
	\subsection{The two research questions and the graph family}
	\label{sec:setup-rqs}
	
	The experiments answer two questions with one design:
	\textbf{(RQ1, which representation?)} does a richer code/semantic structure improve
	commit-level prediction, and which structural representation is best?
	\textbf{(RQ2, which inference method?)} among typical graph-inference paradigms, which
	best exploits the graph?
	
	Both are answered by a single family of graphs of increasing richness, all sharing the
	Core relational layer, evaluated with the fixed five methods:
	\emph{Core}; \emph{Core$+$AST}; \emph{Core$+$CFG}; \emph{Core$+$DFG};
	\emph{Core$+$PDG}; and the deployed final \emph{Core$+$AST$+$CSTG}. Reading a row
	(a method across graphs) answers RQ1 for that method; reading a column (all methods on
	one graph) answers RQ2 for that graph. A structural layer influences prediction through
	exactly one channel --- the per-commit change-tokens \code{\{edge\}:\{type\}} that
	become typed hubs the walks and embeddings traverse --- so swapping the representation
	means swapping the hub stream while everything else is held fixed. On the final graph,
	the hubs additionally include CSTG \node{Term} nodes via \edge{MENTIONS}, so the same
	five methods traverse the semantic signal too. (Token-seq is included in the
	representation \emph{ablation} but is not part of the final six-graph family, being a
	deliberately weak rival.)
	
	\subsection{Online evaluation protocol}
	\label{sec:setup-protocol}
	
	\paragraph{Prequential, predict-then-grow.} Commits are streamed in chronological
	order (\code{author\_ts}). The protocol is \emph{prequential}: for each commit we
	(i)~predict its label from a model and graph state built only from strictly earlier
	commits, (ii)~reveal the true label, (iii)~grow the graph state (add the commit's
	delta edges, term hubs, recency counters, co-occurrence counts), and (iv)~periodically
	re-learn on the expanded window. This exactly reproduces the information available at
	deployment time and structurally precludes look-ahead leakage.
	
	\paragraph{Warm-up, blocks, and refit schedule.} The first \textbf{$40\%$} of the
	stream ($W = \lfloor 0.40\times 8{,}059\rfloor = 3{,}223$ commits) is a warm-up used
	only to initialise models and is \emph{not} scored; the remaining $4{,}836$ commits are
	the evaluation stream. Predictions are made per commit but models are refreshed on a
	\textbf{block} granularity of $200$ commits (the expanding-window retraining that lets
	the models adapt to concept drift). This warm-up/block schedule, the classifier family,
	and all hyper-parameters are held \emph{identical} across every representation, method,
	and fusion so that the only independent variable in each comparison is the one under
	study.
	
	\paragraph{Leakage controls (stated explicitly).} All scalers, vectorisers, IDF
	weights, NPMI graphs, term-risk priors, embedding bases, and stacking regressions are
	fit on past data only; PPR class seeds never include the commit under test; the
	graph-of-words weighting and recency lookups are commit-local; and the
	\code{HashingVectorizer} used for the CSTG text stream is stateless (a fixed hash with
	no learned vocabulary), so it cannot transport information from a future commit into an
	earlier row. The single-split controlled comparison used against the external text
	baseline (\secref{sec:setup-baselines}) is clearly separated from the deployed online
	path and never mixed with it.
	
	\paragraph{Online operating point.} Because the task is imbalanced and drifting, the
	hard-decision threshold is itself tuned \emph{online}: it is re-tuned on the past-only
	predictions every $150$ commits to maximise buggy-F1 on already-seen data, giving a
	leakage-free deployment operating point at which all threshold-based metrics are read.
	Rolling trajectories are computed over a trailing window of $\text{ROLL}=800$ commits
	(a visualisation smoothing only --- it never enters the model or the reported scalars,
	which are computed once over the entire evaluation stream).
	
	\subsection{Evaluation metrics}
	\label{sec:setup-metrics}
	
	Because the positive class is rare and drifting, no single metric suffices; we report a
	\textbf{seven-metric} suite plus two effort-aware measures, and treat the
	threshold-free and class-balanced metrics as primary.
	
	\begin{itemize}
		\item \textbf{Threshold-free ranking:} \emph{ROC-AUC} (a.k.a.\ AUC) and, where
		reported, \emph{PR-AUC} (area under precision--recall; its random baseline equals the
		positive rate, which is why absolute PR-AUC falls in the low-defect tail even for a
		good model).
		\item \textbf{Operating-point metrics at the online-tuned threshold:} \emph{Precision},
		\emph{Recall} (buggy-class sensitivity), \emph{Buggy-F1} (F1 of the positive class),
		\emph{Macro-F1} (unweighted mean of both classes' F1 --- robust to imbalance),
		\emph{G-Mean} ($\sqrt{\text{TPR}\cdot\text{TNR}}$), and \emph{MCC} (Matthews
		correlation --- the single most reliable summary under imbalance).
		\item \textbf{Effort-aware (deployment triage) metrics:} normalised
		$P_{\mathrm{opt}}$ (Kamei et al.: how close the model's inspection order is to the
		ideal buggy-first, least-effort-first order, normalised so $0.5$ is random and $1.0$
		optimal) and \emph{ACC@20\%LOC} (recall achieved when only the top $20\%$ of churned
		lines are inspected). Effort is measured as \code{la}$+$\code{ld}.
	\end{itemize}
	The two \emph{class-balanced} metrics --- Macro-F1 and G-Mean --- are the ones we
	prioritise for model selection, precisely because they do not let a model win by
	exploiting the majority class in the low-defect regime. We emphasise that, given the
	concept drift, the meaningful reading of an online trajectory is \emph{relative} (which
	method/graph stays highest) rather than the absolute end-of-stream value, since the
	falling positive rate depresses every method's tail together.
	
	\subsection{Fusion selection rule}
	\label{sec:setup-fusion-selection}
	
	The deployed fusion is not hand-picked. We evaluate \emph{all} $2^5-1=31$ combinations
	of the five methods on the final graph, score each on the full seven-metric suite, and
	select the smallest combination whose balanced score
	$(\text{Macro-F1}+\text{G-Mean})/2$ is within a fixed tolerance ($0.005$) of the best ---
	a transparent trade-off of \emph{load} (fewer methods, cheaper online) against
	\emph{performance}. This selects $F=\text{RN}{+}\text{PPR}$. We then test, on top of
	$F$, whether the explicit CSTG channel $G$ and/or the non-graph JIT metrics $M$ add
	anything: $G$ helps markedly and $M$ does not (adding $M$ alone hurts), so the deployed
	model is $F{+}G=\text{RN}{+}\text{PPR}{+}\text{CSTG}$. Reporting the full combinatorial
	ablation, and choosing by an explicit rule, is what makes the deployed model
	defensible rather than cherry-picked.
	
	\subsection{Baselines and comparison protocols}
	\label{sec:setup-baselines}
	
	We compare against baselines at two levels.
	\textbf{(i) Within-project JIT baselines:} logistic regression, random forest, and
	gradient boosting over the twelve ApacheJIT metrics --- the standard JIT-defect
	predictor --- evaluated under the identical online protocol, plus a naive prior-rate
	predictor. \textbf{(ii) A strong external text baseline:} a $768$-dimensional
	diff-text TF-IDF model with a random forest, the strongest published lexical baseline.
	Because that baseline was reported on a single chronological split (train first $70\%$,
	test last $15\%$; a hard, low-defect tail with $\approx\!7.5\%$ positives), we
	reproduce that \emph{exact} split and test set for a like-for-like comparison, using a
	matched classifier so any difference comes from the features and not the learner; these
	single-split numbers are reported separately from, and never conflated with, the
	fully-online results. As a sanity check our JIT-metrics model reproduces the baseline
	logistic regression on that split to three decimals, confirming the protocols are
	aligned. For the CSTG internal ablation specifically, we additionally report a
	\emph{same-learner} (matched random forest) plain-TF-IDF row so the ``does the CSTG
	representation beat a bag-of-words'' comparison is under one learner rather than mixing
	linear and non-linear models.
	
	\subsection{Scalability, complexity, and statistical-significance setup}
	\label{sec:setup-scalability}
	
	Beyond accuracy, the paper characterises the system's cost and the statistical
	reliability of its comparisons. All of this is computed \emph{without rebuilding} the
	knowledge graph: final sizes and growth come from read-only queries against the built
	graph and from attributing delta edges to their source commits in online order;
	prediction latency is measured by replaying the cached graph family in memory; and
	per-commit build/modify cost is obtained by an instrumented replay of the real
	parse-and-diff work on a stratified sample of commits, cross-checked against an
	analytical operation-count proxy.
	
	\begin{itemize}
		\item \textbf{Statistical profile.} Per-layer node/edge counts, type and token
		vocabularies, per-file size distributions (with Gini and skew), and a discrete
		power-law fit of the per-commit change-size distribution (heavy-tailed, exponent
		$2<\alpha<3$ on every layer).
		\item \textbf{Per-commit growth and modification.} Cumulative nodes/edges over the
		commit stream and over calendar time (overlaid on the drifting bug rate), the
		per-commit change-size distribution, and the buggy-vs-benign change-size separation
		--- which is large and highly significant (Mann--Whitney $p\ll 10^{-100}$ on every
		layer; buggy commits make on the order of $2$--$3\times$ larger structural changes),
		the basic reason the graph is predictive.
		\item \textbf{Time and space complexity.} Build/modify cost is shown to be near-linear
		in file size ($R^2\approx0.99$), i.e.\ per-commit work is $O(\text{change})$ and space
		is $O(\text{alive nodes}+\text{retained history})$; prediction latency scales with
		graph density, and the deployed graph-native fusion predicts at sub-millisecond per
		commit (thousands of commits/second), so the whole pipeline is comfortably real-time
		and GPU-free. Each per-commit method's cost is annotated with its known asymptotic
		complexity and validated against the measured block-time-vs-graph-size growth.
		\item \textbf{Statistical significance.} Headline comparisons are tested on paired
		per-commit out-of-sample scores: DeLong's test for ROC-AUC differences (Holm-corrected
		within each family of comparisons), paired bootstrap tests and $95\%$ confidence
		intervals for PR-AUC/MCC/G-Mean, and Cliff's~$\delta$ effect sizes. This lets us state
		which accuracy gaps are significant (e.g.\ graph richness significantly helps the
		informative methods; PPR is significantly the best method on the final graph; the
		semantic channel $F{+}G$ significantly improves every headline metric over $F$) and
		which are within noise.
	\end{itemize}
	
	\subsection{Implementation and reproducibility}
	\label{sec:setup-impl}
	
	The graph store is Neo4j (final on-disk footprint $\approx\!3.3$\,GB; working footprint
	$\approx\!5.6$\,GB with transaction logs). Java is parsed with a single deterministic
	parser (\texttt{javalang}) for the whole pipeline, in killable subprocesses with a
	watchdog so a pathological file cannot stall the build; a persistent parser worker and
	an in-memory before-graph cache keep the online build tractable, and a range index on
	each node label's \code{id} keeps per-commit cost constant as the graph grows (an
	absent index was, at one point, the dominant cost --- turning an $O(\text{change})$
	build into an accidentally quadratic one). All inference code is CPU-only Python
	(scikit-learn linear models, custom sparse graph routines, gensim/PyKEEN for the
	embeddings where used); no GNN and no GPU are required. Feature streams and per-method
	scores are cached to disk so every experiment, table, figure, and statistical test can
	be regenerated without a live database or a graph rebuild, and companion Jupyter
	notebooks reproduce all figures from those caches. Randomised components (embedding
	initialisation, bootstrap resampling) use fixed seeds. Determinism of node identity
	across versions (the same source always yields the same parse-local ids) is what makes
	the online build itself reproducible.
	
	\paragraph{Summary of the final configuration.} The deployed predictor is a
	fully-online, GPU-free, multimodal commit knowledge graph on \texttt{apache/groovy}
	($8{,}059$ ApacheJIT commits, $\approx\!20\%$ buggy, 2003--2019): a Core process layer,
	an AST delta layer, and the CSTG semantic-text layer, over which the fusion
	$F{+}G=\text{RN}{+}\text{PPR}{+}\text{CSTG}$ predicts each commit prequentially from
	strictly-past state, warm-up $40\%$ ($3{,}223$ commits), block $200$, evaluated on the
	remaining $4{,}836$ commits with a seven-metric plus effort-aware suite and paired
	significance testing.

	% =====================================================================
	\section{Environment, Software Versions, and Model-Configuration Checklist}
	\label{sec:setup-checklist}
	% =====================================================================
	For completeness and reproducibility, this section answers a standard checklist of
	setup details --- operating system and hardware, software and library versions,
	training hyper-parameters, model architecture, and stopping/tuning. Several checklist
	items are framed for \emph{deep-learning} pipelines (optimiser, learning rate, batch
	size, epochs, hidden-layer dimensions, early stopping). Our final methodology is
	deliberately \textbf{GPU-free and GNN-free}: the inference models are classical
	graph algorithms (random-walk / label-propagation / matrix-factorisation) with
	\emph{linear} (logistic-regression) heads, plus one compact bilinear knowledge-graph
	embedding trained by a short custom SGD. Where a checklist item has no analogue in
	this design, we say so explicitly and explain why, rather than inventing a value.

	\subsection{Operating system and hardware}
	\label{sec:setup-hw}
	All building and all experiments were run on a single commodity laptop, with no GPU
	and no cluster:
	\begin{itemize}[nosep,leftmargin=1.4em]
		\item \textbf{Operating system:} Microsoft Windows~11 (64-bit, build $26{,}200$).
		The pipeline is OS-agnostic (pure Python plus a Dockerised Neo4j) and also runs on
		Linux/macOS; the Neo4j administration helpers assume a Docker-hosted server.
		\item \textbf{CPU:} Intel Core i7-13620H (13th generation), $10$ physical cores /
		$16$ logical threads, $2.4$\,GHz base clock. \textbf{No GPU is used at any stage}
		(build, inference, or analysis).
		\item \textbf{Memory (RAM):} $16$\,GB. The whole pipeline is designed to run within
		this budget: the online build stores only one evolving structure per file (not a
		snapshot per commit), and the inference/analysis code streams and caches to disk, so
		no stage requires loading the full graph into memory at once.
		\item \textbf{Graph store:} Neo4j runs in a Docker container; the built Groovy graph
		occupies $\approx\!3.3$\,GB on disk ($\approx\!5.6$\,GB with transaction logs).
	\end{itemize}

	\subsection{Software and library versions}
	\label{sec:setup-versions}
	The exact versions used are listed in Table~\ref{tab:setup-versions}. The final
	pipeline depends only on the mainstream scientific-Python stack, a Java parser, and a
	Neo4j driver; there is \textbf{no deep-learning framework} (no PyTorch, TensorFlow, or
	Keras) anywhere on the critical path, consistent with the CPU-only design.

	\begin{table}[t]
		\centering
		\caption{Software and library versions used for all builds and experiments.}
		\label{tab:setup-versions}
		\begin{tabular}{lll}
			\toprule
			Component & Version & Role \\
			\midrule
			Python            & 3.13.5   & runtime for the entire pipeline \\
			Neo4j             & 2026.05 (Community) & graph store (Dockerised) \\
			\code{neo4j} (driver) & 6.2.0 & Python$\to$Neo4j Bolt driver \\
			git               & 2.45.1   & reads per-commit file blobs / diffs \\
			javalang          & 0.13.0   & Java parser (AST/CFG/DFG/PDG/SEQ builders) \\
			NumPy             & 2.2.6    & arrays; the graph-method kernels \\
			SciPy             & 1.16.2   & sparse matrices; statistics (Mann--Whitney, ranks) \\
			scikit-learn      & 1.7.2    & logistic-regression heads, TF-IDF, RF/GB baselines \\
			pandas            & 2.3.2    & CSV / tabular handling \\
			networkx          & 3.5      & in-memory graph construction in the builders \\
			matplotlib        & 3.10.6   & figures \\
			nbformat / nbconvert & 5.10.4 / 7.17.1 & build and execute the result notebooks \\
			PyYAML            & 6.0.2    & per-project configuration (reproducibility package) \\
			\bottomrule
		\end{tabular}
	\end{table}

	\paragraph{A note on gensim / PyKEEN.} Earlier iterations experimented with
	\code{gensim} (node2vec/DeepWalk) and \code{PyKEEN} (TransE/DistMult). These are
	\textbf{not} used by the final methodology: the deployed DeepWalk embedding (DW) is a
	self-contained truncated-SVD of a PPMI matrix, and the deployed knowledge-graph
	embedding (KGE) is our own compact DistMult in NumPy (below). No external embedding
	library is on the final critical path.

	\subsection{Training hyper-parameters}
	\label{sec:setup-hparams}
	The checklist items \emph{optimiser}, \emph{learning rate}, \emph{batch size}, and
	\emph{number of epochs} are properties of gradient-descent neural training. Our final
	models are not neural, so most of these \textbf{do not apply}; the few genuine
	training knobs are given precisely.

	\begin{description}[leftmargin=1.4em]
		\item[Optimiser.] \emph{Mostly not applicable.} The linear heads on every graph
		method and on the $M$ and $G$ fusion channels are \textbf{balanced logistic
		regressions solved by L-BFGS} (scikit-learn, \code{solver="lbfgs"},
		\code{class\_weight="balanced"}, \code{max\_iter}$=1000$--$2000$); L-BFGS is a
		quasi-Newton batch solver, so there is no stochastic optimiser, no learning-rate
		schedule, and no mini-batching. The \emph{only} stochastic-gradient component in the
		whole system is the DistMult KGE, trained with plain SGD.
		\item[Learning rate.] \emph{Only the KGE has one:} the DistMult embedding uses a
		fixed SGD learning rate of $\eta=0.05$. All logistic-regression heads are L-BFGS and
		have no learning rate. DeepWalk (DW) is closed-form (SVD) and has none.
		\item[Batch size.] \emph{Not applicable in the neural sense.} There is no
		mini-batch training. The one gradient method (KGE) processes its past triples in
		internal chunks of $2{,}048$ for vectorised updates, but this is an implementation
		detail of a full-pass SGD, not a tuned batch-size hyper-parameter. The evaluation
		protocol does have a \emph{block} granularity ($200$ commits) at which models are
		refreshed, but that is a streaming-protocol constant (\secref{sec:setup-protocol}),
		not a training batch size.
		\item[Number of epochs.] \emph{Only the KGE has epochs:} the DistMult SGD runs
		$3$ epochs over the strictly-past triples at each refit; negative sampling uses
		$3$ corrupted tails per positive. No other model has an epoch count: the
		logistic heads run to L-BFGS convergence (capped at \code{max\_iter}), and DW/RN/PPR/LP
		are single-pass or fixed-iteration (below), not epoch-based.
	\end{description}

	\noindent The fixed-iteration counts of the non-embedding graph methods (which play the
	role ``epochs'' would in a neural model) are: \textbf{PPR} --- $40$ power-iterations of
	random-walk-with-restart, restart probability $\alpha=0.15$, run once per class seed set
	per block; \textbf{LP} --- $3$ propagation steps with clamping on known past labels;
	\textbf{RN} --- a single one-hop weighted vote (no iteration). These are held fixed
	across all projects and all graphs.

	\subsection{Model architecture}
	\label{sec:setup-arch}
	The checklist item \emph{hidden-layer dimensions} presupposes a multi-layer neural
	network. Our models have \textbf{no hidden layers} --- there is no multi-layer
	network, and specifically \textbf{no GNN} (a heterogeneous GNN is explicitly scoped as
	future work, not part of this study). The only learned latent dimensionalities are the
	two shallow embedding sizes already reported: \textbf{DeepWalk (DW) $=64$} and
	\textbf{DistMult KGE $=32$}. Beyond those:
	\begin{itemize}[nosep,leftmargin=1.4em]
		\item \textbf{RN, PPR, LP} produce a scalar risk score directly from the
		commit--hub incidence / transition matrix; they have no latent layer at all.
		\item \textbf{DW} embeds commits into $\mathbb{R}^{64}$ by truncated-SVD of the PPMI
		of the commit--hub co-occurrence (Levy--Goldberg form), then a logistic head; there
		is no hidden layer, only this one factorisation dimension.
		\item \textbf{KGE (DistMult)} embeds hubs and relations into $\mathbb{R}^{32}$ with a
		bilinear score; commits are folded in as the mean of their hubs' embeddings, then a
		logistic head. Again a single latent dimension, no hidden layers.
		\item \textbf{Text/structural feature widths} (not ``hidden'' layers, but the input
		dimensionalities worth recording): the structural TF--IDF and the CSTG text stream
		are hashed to $2^{18}$ features (a fixed, stateless hashing vectoriser); the CSTG
		typed-mass vector is $5$-dimensional; the JIT-metric channel is $12$-dimensional.
	\end{itemize}
	So the ``architecture'' of the deployed model is: five shallow graph-inference scorers
	(two of them with a $64$- or $32$-dim embedding) whose scalar outputs are combined by a
	single balanced logistic-regression stacker --- no hidden layers, no GNN.

	\subsection{Stopping and tuning}
	\label{sec:setup-stopping}
	\begin{description}[leftmargin=1.4em]
		\item[Early-stopping criteria.] \emph{Not used, and not applicable.} Because the
		models are non-neural, there is no validation-loss-based early stopping. The KGE
		SGD runs a fixed $3$ epochs; the logistic heads stop at L-BFGS convergence (or the
		\code{max\_iter} cap); the fixed-iteration methods (PPR $40$, LP $3$) run to their
		set iteration count. The streaming protocol itself has no ``stopping'': every
		post-warm-up commit is scored, and the run ends only when the commit stream ends.
		\item[Hyper-parameter tuning strategy.] We deliberately \textbf{do not tune the
		model hyper-parameters per project or per graph}, and this is a design choice, not
		an omission: the whole study is a \emph{controlled comparison}, so the protocol
		constants (warm-up $40\%$, block $200$, roll $800$, refit every $5$ blocks), the
		embedding sizes ($64$/$32$), the iteration counts (PPR $40$, LP $3$, KGE
		$3$ epochs / $\eta=0.05$ / $3$ negatives), and the classifier family (balanced L-BFGS
		logistic regression) are \textbf{held fixed across every representation, method,
		fusion, and project}. Fixing them is what makes cross-configuration and
		cross-project differences attributable to the factor under study rather than to
		per-cell tuning. The one thing that \emph{is} selected --- the deployed fusion
		formula --- is chosen not by a hyper-parameter search but by the transparent,
		reported combinatorial rule of \secref{sec:setup-fusion-selection} (evaluate all
		$2^5-1$ method combinations, take the smallest within a fixed tolerance of the best
		balanced score). The threshold at which hard decisions are read is likewise not a
		tuned constant but is re-derived \emph{online} from past-only data every $150$
		commits (\secref{sec:setup-protocol}). The default hyper-parameter values were set
		once, up front, from standard choices in the literature (e.g.\ $\alpha=0.15$ for
		PPR) and small sanity checks; there was no automated grid/random/Bayesian search
		over them, precisely to avoid over-fitting the evaluation.
	\end{description}

	\paragraph{Cross-project note.} Every constant above is identical across all subject
	projects (Groovy and the additional Apache projects processed with the reproducibility
	package); only the data changes. The single genuinely per-project quantity is the base
	commit that seeds the structural layers (\secref{sec:setup-data}), which is a property
	of each repository's history, not a tuned hyper-parameter.

\end{document} from the parent .tex.
% =====================================================================
\documentclass[11pt,a4paper]{article}
\usepackage[margin=2.4cm]{geometry}
\usepackage[T1]{fontenc}
\usepackage[utf8]{inputenc}
\usepackage{lmodern}
\usepackage{amsmath,amssymb}
\usepackage{booktabs}
\usepackage{array}
\usepackage{enumitem}
\usepackage{xcolor}
\usepackage[hidelinks]{hyperref}

% ---- macros used throughout (match the main document's conventions) ----
\providecommand{\code}[1]{\texttt{#1}}
\providecommand{\node}[1]{\textsf{#1}}
\providecommand{\edge}[1]{\textsc{#1}}

% ---- helper for safe section references (avoids any parsing issues with \S\ref) ----
\newcommand{\secref}[1]{\S\ref{#1}}

\begin{document}
	
	\section{Experimental Setup}
	\label{sec:experimental-setup}
	
	This section specifies, in full, the data, the knowledge-graph construction, the
	inference models, the online evaluation protocol, the metrics, the baselines, and
	the reproducibility conditions under which every result in this paper was produced.
	Two principles govern the whole design and are worth stating up front because they
	recur throughout. \textbf{First, everything is \emph{just-in-time (JIT) and fully
			online}:} the graph grows one commit at a time in true chronological order, and a
	commit is always predicted from a model and a graph state built \emph{only} from
	strictly earlier commits --- there is no random split, no look-ahead, and no
	statistic is ever fit on data that includes or postdates the commit under test.
	\textbf{Second, the study is designed as a controlled comparison:} when we ask which
	code representation or which inference method is best, exactly one factor is varied
	at a time while the protocol, the classifier family, the warm-up, and the block
	schedule are held fixed, so any measured difference is attributable to that factor
	and not to a confound.
	
	\subsection{Subject program, commit scope, and labels}
	\label{sec:setup-data}
	
	\paragraph{Project.} All experiments are conducted on \texttt{apache/groovy}, a
	mature, long-lived Java project (the Groovy programming language and its compiler,
	runtime, and tooling). We work from a full local \texttt{git} clone, so the raw
	material for every commit --- its message, its parent(s), and the exact byte content
	of every file before and after the change --- is available offline and deterministic.
	
	\paragraph{Labelled commit scope.} The study set is the \textbf{$8{,}059$ commits}
	of \texttt{apache/groovy} that appear in the \textbf{ApacheJIT} dataset
	(\texttt{data/apachejit/projects/apache\_groovy.csv}), spanning \textbf{2003--2019}.
	ApacheJIT provides, for each commit, a \emph{buggy/benign} label derived by an
	SZZ-style procedure (a commit is buggy if a later fix commit points back to it) and
	the twelve standard JIT change metrics. These $8{,}059$ commits are the only ones we
	treat as the labelled study set and the only ones grown into the AST/structural and
	semantic layers. The surrounding history is retained in the database as
	\node{Commit} nodes for context (the full clone has $16{,}887$ commits), but only the
	ApacheJIT commits are flagged \code{in\_jit=true} and are the units of prediction.
	
	\paragraph{Class balance and its consequences.} The labelled set is imbalanced:
	\textbf{$1{,}614$ buggy / $6{,}445$ benign}, i.e.\ a $\approx\!20\%$ positive rate
	overall. Crucially, the positive rate is \emph{not stationary}: it rises from
	$\approx\!6\%$ in 2003 to a peak of $\approx\!39\%$ in 2012 and then falls back to
	$\approx\!6\%$ by 2019 (measured directly from the labels). This pronounced
	\emph{concept drift} is the single most important property of the data for our
	purposes: it is why a random or single-split evaluation is inappropriate, why we
	adopt a time-aware online protocol, and why the tail of every online trajectory
	declines (the achievable precision/recall shrink as positives become rare). We
	return to this when defining the protocol (\secref{sec:setup-protocol}) and the
	metrics (\secref{sec:setup-metrics}).
	
	\paragraph{Base snapshot.} The structural layers are seeded once at a fixed
	\emph{base} commit --- \texttt{408b29851d7bbe4d343340832297e4be7e0c5578}, the first
	substantial Java commit --- whose $120$ \texttt{.java} files (${\approx}46$k AST
	nodes) constitute the initial code structure onto which all subsequent per-commit
	changes are laid. The base is a property of the \emph{code}, not of the labelled
	set; it merely fixes the starting state of the evolving graph.
	
	\subsection{Knowledge-graph construction and online growth}
	\label{sec:setup-kg}
	
	The knowledge graph is a \emph{heterogeneous information network} (HIN): typed nodes
	and typed edges in one store (Neo4j). It fuses three families of evidence that the
	defect-prediction literature usually treats separately --- \emph{process} evidence
	(who changed what, when, and how often), \emph{code-structure} evidence (the
	syntactic/semantic shape of the change and its surroundings), and \emph{textual}
	evidence (the natural-language and lexical content of the change). The design
	rationale for a graph (rather than a flat feature table) is that anchoring each
	commit's change to \emph{shared, persistent} code entities is what makes relational
	and neighbourhood reasoning across the entire history possible.
	
	\paragraph{Core (process / relational) layer.} The Core layer is the classical
	commit graph and is present in \emph{every} experimental configuration. Its node
	types are \node{Project}, \node{Branch}, \node{Developer}, \node{Commit},
	\node{File}, and \node{Issue}; its edges encode authorship
	(\node{Commit}$\xrightarrow{\text{\edge{AUTHORED\_BY}}}$\node{Developer}), the
	history spine (\node{Commit}$\xrightarrow{\text{\edge{PARENT\_OF}}}$\node{Commit}),
	file changes
	(\edge{ADDED}/\edge{MODIFIED}/\edge{DELETED}/\edge{RENAMED\_FROM}/\edge{RENAMED\_TO}
	between \node{Commit} and \node{File}), and issue links
	(\node{Commit}$\xrightarrow{\text{\edge{FIXES\_ISSUE}}}$\node{Issue}). Every
	\node{Commit} also carries the twelve ApacheJIT change metrics (size \code{la},
	\code{ld}; dispersion \code{nf}, \code{nd}, \code{ns}, \code{ent}; history and
	experience \code{age}, \code{ndev}, \code{nuc}, \code{aexp}, \code{arexp},
	\code{asexp}) and its \code{author\_ts} (the UNIX timestamp that defines online
	arrival order), \code{buggy} label, and \code{jit\_year}.
	
	\paragraph{Structural (code) layers as \emph{delta} graphs.} On top of the Core, we
	materialise the code structure of every changed \texttt{.java} file. The key modelling
	decision --- and a central contribution --- is that a commit is not attached to a
	per-commit snapshot blob; it is attached, by four typed \emph{delta} edges
	(\edge{ADDS}, \edge{REMOVES}, \edge{UPDATES}, \edge{MOVES}), directly to the specific
	code nodes it changed, laid on top of a single evolving structure per file. Concretely,
	for each tracked file the graph holds one \emph{current} structure; when a commit
	touches the file we (i) parse the before- and after-states with the \emph{same}
	deterministic builder, (ii) diff them, (iii) write the delta edges from the
	\node{Commit} to the affected nodes, and (iv) evolve the stored structure so the next
	commit diffs against the correct state. Node identity is threaded across versions by a
	deterministic, per-file parse-local-id map, so the ``after'' map of commit $N$ is
	byte-identical to the ``before'' map of the next commit that touches the file when the
	intervening source is unchanged. Removed nodes are \emph{retained} and flagged
	(\code{alive=false}, \code{removed\_by}), so the complete change history is queryable
	while the live structure always reflects the current code. This makes storage grow
	with the \emph{amount of change}, not with history-length $\times$ file-size, and it is
	what makes the whole graph online-extensible.
	
	We instantiate this delta-graph scheme for \textbf{five} structural representations of
	increasing coarseness, all built per changed \texttt{.java} file and attached at the
	\node{File} level:
	\begin{itemize}
		\item \textbf{AST} --- the abstract syntax tree; one node per syntactic construct
		(declaration, statement, expression, operator, literal, and each identifier/operator
		leaf). Finest granularity ($136$ node types).
		\item \textbf{CFG} --- the intra-procedural control-flow graph; one node per statement
		or control construct ($18$ node types).
		\item \textbf{DFG} --- the def--use data-flow graph over local variables; one node per
		definition/use site ($7$ node types).
		\item \textbf{PDG} --- the program-dependence graph; the CFG's statement vertices plus
		control-dependence and data-dependence \emph{edges} (no new nodes; $18$ node types).
		We label it PDG rather than CPG because we deliberately keep dependence at statement
		granularity and do \emph{not} embed the AST subtree under each statement; a canonical
		code property graph would contain the AST as a subgraph and is out of scope.
		\item \textbf{Token-seq} --- a deliberately weak flat statement sequence (statement
		nodes chained by \edge{NEXT}, no control or data structure); a lower-bound rival used
		only in the representation comparison.
	\end{itemize}
	Table~\ref{tab:setup-layers} reports the materialised scale of each layer. The
	${\approx}10\times$ gap between AST ($3.70$\,M nodes) and the others ($0.26$--$0.47$\,M)
	is a pure \emph{granularity} effect: a statement such as \code{int x = foo(a,b)+1;} is
	one node in CFG/PDG but a dozen-plus nodes in the AST. CFG and PDG have identical node
	counts because the PDG adds only dependence edges to the CFG's vertex set.
	
	\begin{table}[t]
		\centering
		\caption{Materialised scale of each structural layer and its per-commit change-token
			stream (Groovy, $8{,}059$ commits). ``Token types'' counts the distinct
			\code{\{edge\}:\{node-type\}} change-tokens; ``Commits w/ tok.'' is the number of
			labelled commits that emit at least one token in that layer.}
		\label{tab:setup-layers}
		\begin{tabular}{lrrrrr}
			\toprule
			Layer & Nodes & Delta-edges & Node types & Token types & Commits w/ tok. \\
			\midrule
			AST       & 3{,}699{,}439 & 3{,}614{,}922 & 136 & 314 & 7{,}523 \\
			CFG       &   372{,}488 &   408{,}301 & 18 & 53 & 6{,}194 \\
			DFG       &   469{,}662 &   648{,}553 &  7 & 27 & 6{,}066 \\
			PDG       &   372{,}488 &   411{,}869 & 18 & 63 & 6{,}261 \\
			Token-seq &   257{,}708 &   293{,}287 & 16 & 47 & 6{,}011 \\
			\bottomrule
		\end{tabular}
	\end{table}
	
	\paragraph{The differ, and why matching is fair across layers.} The AST layer is
	diffed as a \emph{tree} by bottom-up subtree hashing (identical subtrees in the before
	and after trees are matched greedily, preferring equal source positions; unmatched
	before-nodes are deletions, unmatched after-nodes insertions, matched-but-relabelled
	nodes updates, reparented nodes moves). The four non-tree layers (CFG/DFG/PDG/Token-seq)
	are general typed graphs and are diffed by a single \emph{node-identity-lite} matcher
	applied \emph{identically} to each: matching is scoped within a method (\code{class::name},
	stable across commits) and keyed by \emph{ordinal} (the $k$-th node of a given type in a
	method maps to the $k$-th such node in the other version), which is robust to line
	shifts. Using one shared differ for the alternatives is the fairness guarantee for the
	representation comparison: only the subgraph varies, never the difference model. We
	disclose that the AST uses a tree differ while the others use the ordinal matcher --- a
	deliberate specialisation (trees vs.\ general graphs), not a hidden confound.
	
	\paragraph{Semantic-text layer (CSTG).} The third first-class layer is the
	\textbf{Commit Semantic-Text Graph (CSTG)}, which models the \emph{textual}
	description of a change (commit message $+$ diff text) as a graph grounded in code.
	Its node types are \node{Term} (a typed vocabulary term, with \code{kind} $\in$
	\{\emph{code}-identifier, \emph{bug}-lexicon, \emph{action}, \emph{error}-type,
	\emph{nl}\}) and \node{Intent} (the change-intent class: fix, feat, refactor, \dots);
	its edges are
	\node{Commit}$\xrightarrow{\text{\edge{MENTIONS}\,(w=\text{TW-IDF})}}$\node{Term} (top
	terms per commit, weighted by TextRank centrality in the commit's word-graph $\times$
	IDF), \node{Term}$\xrightarrow{\text{\edge{COOCCURS}\,(\text{npmi})}}$\node{Term} (the
	global normalised-pointwise-mutual-information term graph),
	\node{Term}$\xrightarrow{\text{\edge{GROUNDS\_IN}}}$\node{ASTNode} (a code term linked
	to the alive identifier leaf it names), and
	\node{Commit}$\xrightarrow{\text{\edge{HAS\_INTENT}}}$\node{Intent}. On Groovy the CSTG
	comprises $20{,}377$ \node{Term} nodes, $188{,}435$ \edge{MENTIONS}, $120{,}340$
	\edge{COOCCURS}, $19{,}315$ \edge{GROUNDS\_IN}, and $8{,}059$ \edge{HAS\_INTENT}.
	Theoretically the layer rests on graph-of-words / TW-IDF (Rousseau \& Vazirgiannis),
	the NPMI corpus term graph (cf.\ TextGCN), defect-semantic typing, and term-risk
	propagation over the NPMI graph. Every CSTG statistic that enters prediction is grown
	online, past-only (see \secref{sec:setup-protocol}); the static term graph and
	term-risk values written to Neo4j are a queryable/visualisation snapshot and are
	\emph{not} read by the online predictor.
	
	\paragraph{Data-quality guarantees.} The finalised graph passes a standing integrity
	audit: no file has more than one structure root, removed nodes always carry a
	\edge{REMOVES} edge, and the online-evolved structure reproduces a fresh
	ground-truth parse of the current file state with precision $=$ recall $=1.0$ on the
	validated sample. A large \emph{ghost-AST} inflation ($\approx\!1.58$\,M nodes left by
	package-restructuring renames outside the labelled set) was identified and corrected,
	so the reported node counts reflect the clean final state. Known, documented
	limitations that a reviewer should be aware of: labels are SZZ-derived and carry the
	usual mining noise; non-\texttt{.java} files are not modelled at the code level; and
	a git rename is currently treated as an add-of-new-path in the delta layer.
	
	\subsection{Inference models}
	\label{sec:setup-models}
	
	The target hardware has \textbf{no GPU}; every method is therefore deliberately
	CPU-only and \textbf{GNN-free}, both to be deployable and to constitute strong,
	honest baselines that any future GNN must beat. The final methodology fixes
	\textbf{five canonical graph-inference paradigms}, chosen because they map one-to-one
	onto the recognised families of graph and knowledge-graph inference, each reads a
	\emph{different, complementary} facet of the structure, and --- crucially --- each is a
	graph \emph{algorithm} that runs unchanged on \emph{any} graph in the family, so all
	five are valid rows across every column of the representation study:
	
	\begin{description}
		\item[RN --- Relational neighbour (wvRN).] Collective classification: a commit's
		score is the label-weighted hub-overlap with past commits (one hop through the shared
		file/developer/term hubs). Complexity $O(\mathrm{nnz})$ of a sparse block matmul.
		\item[PPR --- Personalised PageRank (random walk with restart).] The bipartite
		commit--hub graph is seeded from past buggy vs.\ benign commits and a commit is scored
		by its relative proximity to each class. Complexity $O(\text{iters}\cdot|E|)$ per
		block, run once per class seed set.
		\item[LP --- Label propagation.] Spreading activation of past labels through the
		hubs a few steps, clamped on known past labels. Complexity $O(\text{iters}\cdot|E|)$.
		\item[DW --- DeepWalk-style embedding.] A random-walk / matrix-factorisation node
		embedding (PPMI of the commit--hub co-occurrence, truncated-SVD, the Levy--Goldberg
		form) with a logistic head; marginals from past rows only.
		\item[KGE --- Knowledge-graph embedding (DistMult).] A compact bilinear embedding
		over typed $(\text{commit},\text{rel},\text{hub})$ triples restricted to past
		commits, with commits folded in from their hubs' learned embeddings and a logistic
		head.
	\end{description}
	
	The embedding methods (DW, KGE) are refit on the expanding past window every
	\textbf{$5$ blocks}; RN, PPR, and LP are recomputed each block from the strictly-past
	state. Embedding dimensions are $64$ (DW) and $32$ (KGE). All linear heads are
	balanced logistic regressions.
	
	\paragraph{Two non-graph channels used only in fusion.} Two additional signals are
	\emph{not} admissible as rows of the representation study (they are not graph
	algorithms) but are evaluated as fusion channels:
	\begin{itemize}
		\item \textbf{$M$ --- JIT metrics:} a prequential logistic model over the twelve
		ApacheJIT change metrics; the classical JIT-defect predictor and our metric-only
		reference channel.
		\item \textbf{$G$ --- CSTG feature channel:} a prequential logistic model over the
		CSTG's per-commit features --- the graph-of-words (TextRank-weighted, add/remove
		polarity-tagged) term stream, the NPMI-propagated past-only term-risk prior, a
		$5$-dimensional typed-mass vector over the defect lexicon, and the intent/consistency
		features --- enriched with an online recency / bug-cache block (per-file and
		per-developer time-since-last-buggy touch, recently-buggy touched-file count, re-fix
		flag, commit velocity), each computed from strictly-past state.
	\end{itemize}
	The commit-\emph{message} text channel ($X$) that appeared in earlier iterations is
	\textbf{removed entirely} from the final methodology (it is a strict subset of $G$ and
	adds nothing over it); it is disabled throughout and reported only as a historical
	step.
	
	\paragraph{Fusion.} Method and channel scores are combined by \emph{prequential
		logistic-regression stacking}: at each block a balanced logistic regression is fit on
	the strictly-past predicted probabilities of the constituent models and applied to the
	current block, refit on the expanding window. The \emph{deployed} model is selected by
	an explicit load-vs-performance rule (\secref{sec:setup-fusion-selection}) and is the
	purely graph-native \textbf{$F{+}G = \text{RN}{+}\text{PPR}{+}\text{CSTG}$}, i.e.\ the
	two-method graph fusion $F=\text{RN}{+}\text{PPR}$ augmented with the CSTG semantic
	channel $G$, \emph{without} the non-graph JIT metrics $M$.
	
	\subsection{The two research questions and the graph family}
	\label{sec:setup-rqs}
	
	The experiments answer two questions with one design:
	\textbf{(RQ1, which representation?)} does a richer code/semantic structure improve
	commit-level prediction, and which structural representation is best?
	\textbf{(RQ2, which inference method?)} among typical graph-inference paradigms, which
	best exploits the graph?
	
	Both are answered by a single family of graphs of increasing richness, all sharing the
	Core relational layer, evaluated with the fixed five methods:
	\emph{Core}; \emph{Core$+$AST}; \emph{Core$+$CFG}; \emph{Core$+$DFG};
	\emph{Core$+$PDG}; and the deployed final \emph{Core$+$AST$+$CSTG}. Reading a row
	(a method across graphs) answers RQ1 for that method; reading a column (all methods on
	one graph) answers RQ2 for that graph. A structural layer influences prediction through
	exactly one channel --- the per-commit change-tokens \code{\{edge\}:\{type\}} that
	become typed hubs the walks and embeddings traverse --- so swapping the representation
	means swapping the hub stream while everything else is held fixed. On the final graph,
	the hubs additionally include CSTG \node{Term} nodes via \edge{MENTIONS}, so the same
	five methods traverse the semantic signal too. (Token-seq is included in the
	representation \emph{ablation} but is not part of the final six-graph family, being a
	deliberately weak rival.)
	
	\subsection{Online evaluation protocol}
	\label{sec:setup-protocol}
	
	\paragraph{Prequential, predict-then-grow.} Commits are streamed in chronological
	order (\code{author\_ts}). The protocol is \emph{prequential}: for each commit we
	(i)~predict its label from a model and graph state built only from strictly earlier
	commits, (ii)~reveal the true label, (iii)~grow the graph state (add the commit's
	delta edges, term hubs, recency counters, co-occurrence counts), and (iv)~periodically
	re-learn on the expanded window. This exactly reproduces the information available at
	deployment time and structurally precludes look-ahead leakage.
	
	\paragraph{Warm-up, blocks, and refit schedule.} The first \textbf{$40\%$} of the
	stream ($W = \lfloor 0.40\times 8{,}059\rfloor = 3{,}223$ commits) is a warm-up used
	only to initialise models and is \emph{not} scored; the remaining $4{,}836$ commits are
	the evaluation stream. Predictions are made per commit but models are refreshed on a
	\textbf{block} granularity of $200$ commits (the expanding-window retraining that lets
	the models adapt to concept drift). This warm-up/block schedule, the classifier family,
	and all hyper-parameters are held \emph{identical} across every representation, method,
	and fusion so that the only independent variable in each comparison is the one under
	study.
	
	\paragraph{Leakage controls (stated explicitly).} All scalers, vectorisers, IDF
	weights, NPMI graphs, term-risk priors, embedding bases, and stacking regressions are
	fit on past data only; PPR class seeds never include the commit under test; the
	graph-of-words weighting and recency lookups are commit-local; and the
	\code{HashingVectorizer} used for the CSTG text stream is stateless (a fixed hash with
	no learned vocabulary), so it cannot transport information from a future commit into an
	earlier row. The single-split controlled comparison used against the external text
	baseline (\secref{sec:setup-baselines}) is clearly separated from the deployed online
	path and never mixed with it.
	
	\paragraph{Online operating point.} Because the task is imbalanced and drifting, the
	hard-decision threshold is itself tuned \emph{online}: it is re-tuned on the past-only
	predictions every $150$ commits to maximise buggy-F1 on already-seen data, giving a
	leakage-free deployment operating point at which all threshold-based metrics are read.
	Rolling trajectories are computed over a trailing window of $\text{ROLL}=800$ commits
	(a visualisation smoothing only --- it never enters the model or the reported scalars,
	which are computed once over the entire evaluation stream).
	
	\subsection{Evaluation metrics}
	\label{sec:setup-metrics}
	
	Because the positive class is rare and drifting, no single metric suffices; we report a
	\textbf{seven-metric} suite plus two effort-aware measures, and treat the
	threshold-free and class-balanced metrics as primary.
	
	\begin{itemize}
		\item \textbf{Threshold-free ranking:} \emph{ROC-AUC} (a.k.a.\ AUC) and, where
		reported, \emph{PR-AUC} (area under precision--recall; its random baseline equals the
		positive rate, which is why absolute PR-AUC falls in the low-defect tail even for a
		good model).
		\item \textbf{Operating-point metrics at the online-tuned threshold:} \emph{Precision},
		\emph{Recall} (buggy-class sensitivity), \emph{Buggy-F1} (F1 of the positive class),
		\emph{Macro-F1} (unweighted mean of both classes' F1 --- robust to imbalance),
		\emph{G-Mean} ($\sqrt{\text{TPR}\cdot\text{TNR}}$), and \emph{MCC} (Matthews
		correlation --- the single most reliable summary under imbalance).
		\item \textbf{Effort-aware (deployment triage) metrics:} normalised
		$P_{\mathrm{opt}}$ (Kamei et al.: how close the model's inspection order is to the
		ideal buggy-first, least-effort-first order, normalised so $0.5$ is random and $1.0$
		optimal) and \emph{ACC@20\%LOC} (recall achieved when only the top $20\%$ of churned
		lines are inspected). Effort is measured as \code{la}$+$\code{ld}.
	\end{itemize}
	The two \emph{class-balanced} metrics --- Macro-F1 and G-Mean --- are the ones we
	prioritise for model selection, precisely because they do not let a model win by
	exploiting the majority class in the low-defect regime. We emphasise that, given the
	concept drift, the meaningful reading of an online trajectory is \emph{relative} (which
	method/graph stays highest) rather than the absolute end-of-stream value, since the
	falling positive rate depresses every method's tail together.
	
	\subsection{Fusion selection rule}
	\label{sec:setup-fusion-selection}
	
	The deployed fusion is not hand-picked. We evaluate \emph{all} $2^5-1=31$ combinations
	of the five methods on the final graph, score each on the full seven-metric suite, and
	select the smallest combination whose balanced score
	$(\text{Macro-F1}+\text{G-Mean})/2$ is within a fixed tolerance ($0.005$) of the best ---
	a transparent trade-off of \emph{load} (fewer methods, cheaper online) against
	\emph{performance}. This selects $F=\text{RN}{+}\text{PPR}$. We then test, on top of
	$F$, whether the explicit CSTG channel $G$ and/or the non-graph JIT metrics $M$ add
	anything: $G$ helps markedly and $M$ does not (adding $M$ alone hurts), so the deployed
	model is $F{+}G=\text{RN}{+}\text{PPR}{+}\text{CSTG}$. Reporting the full combinatorial
	ablation, and choosing by an explicit rule, is what makes the deployed model
	defensible rather than cherry-picked.
	
	\subsection{Baselines and comparison protocols}
	\label{sec:setup-baselines}
	
	We compare against baselines at two levels.
	\textbf{(i) Within-project JIT baselines:} logistic regression, random forest, and
	gradient boosting over the twelve ApacheJIT metrics --- the standard JIT-defect
	predictor --- evaluated under the identical online protocol, plus a naive prior-rate
	predictor. \textbf{(ii) A strong external text baseline:} a $768$-dimensional
	diff-text TF-IDF model with a random forest, the strongest published lexical baseline.
	Because that baseline was reported on a single chronological split (train first $70\%$,
	test last $15\%$; a hard, low-defect tail with $\approx\!7.5\%$ positives), we
	reproduce that \emph{exact} split and test set for a like-for-like comparison, using a
	matched classifier so any difference comes from the features and not the learner; these
	single-split numbers are reported separately from, and never conflated with, the
	fully-online results. As a sanity check our JIT-metrics model reproduces the baseline
	logistic regression on that split to three decimals, confirming the protocols are
	aligned. For the CSTG internal ablation specifically, we additionally report a
	\emph{same-learner} (matched random forest) plain-TF-IDF row so the ``does the CSTG
	representation beat a bag-of-words'' comparison is under one learner rather than mixing
	linear and non-linear models.
	
	\subsection{Scalability, complexity, and statistical-significance setup}
	\label{sec:setup-scalability}
	
	Beyond accuracy, the paper characterises the system's cost and the statistical
	reliability of its comparisons. All of this is computed \emph{without rebuilding} the
	knowledge graph: final sizes and growth come from read-only queries against the built
	graph and from attributing delta edges to their source commits in online order;
	prediction latency is measured by replaying the cached graph family in memory; and
	per-commit build/modify cost is obtained by an instrumented replay of the real
	parse-and-diff work on a stratified sample of commits, cross-checked against an
	analytical operation-count proxy.
	
	\begin{itemize}
		\item \textbf{Statistical profile.} Per-layer node/edge counts, type and token
		vocabularies, per-file size distributions (with Gini and skew), and a discrete
		power-law fit of the per-commit change-size distribution (heavy-tailed, exponent
		$2<\alpha<3$ on every layer).
		\item \textbf{Per-commit growth and modification.} Cumulative nodes/edges over the
		commit stream and over calendar time (overlaid on the drifting bug rate), the
		per-commit change-size distribution, and the buggy-vs-benign change-size separation
		--- which is large and highly significant (Mann--Whitney $p\ll 10^{-100}$ on every
		layer; buggy commits make on the order of $2$--$3\times$ larger structural changes),
		the basic reason the graph is predictive.
		\item \textbf{Time and space complexity.} Build/modify cost is shown to be near-linear
		in file size ($R^2\approx0.99$), i.e.\ per-commit work is $O(\text{change})$ and space
		is $O(\text{alive nodes}+\text{retained history})$; prediction latency scales with
		graph density, and the deployed graph-native fusion predicts at sub-millisecond per
		commit (thousands of commits/second), so the whole pipeline is comfortably real-time
		and GPU-free. Each per-commit method's cost is annotated with its known asymptotic
		complexity and validated against the measured block-time-vs-graph-size growth.
		\item \textbf{Statistical significance.} Headline comparisons are tested on paired
		per-commit out-of-sample scores: DeLong's test for ROC-AUC differences (Holm-corrected
		within each family of comparisons), paired bootstrap tests and $95\%$ confidence
		intervals for PR-AUC/MCC/G-Mean, and Cliff's~$\delta$ effect sizes. This lets us state
		which accuracy gaps are significant (e.g.\ graph richness significantly helps the
		informative methods; PPR is significantly the best method on the final graph; the
		semantic channel $F{+}G$ significantly improves every headline metric over $F$) and
		which are within noise.
	\end{itemize}
	
	\subsection{Implementation and reproducibility}
	\label{sec:setup-impl}
	
	The graph store is Neo4j (final on-disk footprint $\approx\!3.3$\,GB; working footprint
	$\approx\!5.6$\,GB with transaction logs). Java is parsed with a single deterministic
	parser (\texttt{javalang}) for the whole pipeline, in killable subprocesses with a
	watchdog so a pathological file cannot stall the build; a persistent parser worker and
	an in-memory before-graph cache keep the online build tractable, and a range index on
	each node label's \code{id} keeps per-commit cost constant as the graph grows (an
	absent index was, at one point, the dominant cost --- turning an $O(\text{change})$
	build into an accidentally quadratic one). All inference code is CPU-only Python
	(scikit-learn linear models, custom sparse graph routines, gensim/PyKEEN for the
	embeddings where used); no GNN and no GPU are required. Feature streams and per-method
	scores are cached to disk so every experiment, table, figure, and statistical test can
	be regenerated without a live database or a graph rebuild, and companion Jupyter
	notebooks reproduce all figures from those caches. Randomised components (embedding
	initialisation, bootstrap resampling) use fixed seeds. Determinism of node identity
	across versions (the same source always yields the same parse-local ids) is what makes
	the online build itself reproducible.
	
	\paragraph{Summary of the final configuration.} The deployed predictor is a
	fully-online, GPU-free, multimodal commit knowledge graph on \texttt{apache/groovy}
	($8{,}059$ ApacheJIT commits, $\approx\!20\%$ buggy, 2003--2019): a Core process layer,
	an AST delta layer, and the CSTG semantic-text layer, over which the fusion
	$F{+}G=\text{RN}{+}\text{PPR}{+}\text{CSTG}$ predicts each commit prequentially from
	strictly-past state, warm-up $40\%$ ($3{,}223$ commits), block $200$, evaluated on the
	remaining $4{,}836$ commits with a seven-metric plus effort-aware suite and paired
	significance testing.

	% =====================================================================
	\section{Environment, Software Versions, and Model-Configuration Checklist}
	\label{sec:setup-checklist}
	% =====================================================================
	For completeness and reproducibility, this section answers a standard checklist of
	setup details --- operating system and hardware, software and library versions,
	training hyper-parameters, model architecture, and stopping/tuning. Several checklist
	items are framed for \emph{deep-learning} pipelines (optimiser, learning rate, batch
	size, epochs, hidden-layer dimensions, early stopping). Our final methodology is
	deliberately \textbf{GPU-free and GNN-free}: the inference models are classical
	graph algorithms (random-walk / label-propagation / matrix-factorisation) with
	\emph{linear} (logistic-regression) heads, plus one compact bilinear knowledge-graph
	embedding trained by a short custom SGD. Where a checklist item has no analogue in
	this design, we say so explicitly and explain why, rather than inventing a value.

	\subsection{Operating system and hardware}
	\label{sec:setup-hw}
	All building and all experiments were run on a single commodity laptop, with no GPU
	and no cluster:
	\begin{itemize}[nosep,leftmargin=1.4em]
		\item \textbf{Operating system:} Microsoft Windows~11 (64-bit, build $26{,}200$).
		The pipeline is OS-agnostic (pure Python plus a Dockerised Neo4j) and also runs on
		Linux/macOS; the Neo4j administration helpers assume a Docker-hosted server.
		\item \textbf{CPU:} Intel Core i7-13620H (13th generation), $10$ physical cores /
		$16$ logical threads, $2.4$\,GHz base clock. \textbf{No GPU is used at any stage}
		(build, inference, or analysis).
		\item \textbf{Memory (RAM):} $16$\,GB. The whole pipeline is designed to run within
		this budget: the online build stores only one evolving structure per file (not a
		snapshot per commit), and the inference/analysis code streams and caches to disk, so
		no stage requires loading the full graph into memory at once.
		\item \textbf{Graph store:} Neo4j runs in a Docker container; the built Groovy graph
		occupies $\approx\!3.3$\,GB on disk ($\approx\!5.6$\,GB with transaction logs).
	\end{itemize}

	\subsection{Software and library versions}
	\label{sec:setup-versions}
	The exact versions used are listed in Table~\ref{tab:setup-versions}. The final
	pipeline depends only on the mainstream scientific-Python stack, a Java parser, and a
	Neo4j driver; there is \textbf{no deep-learning framework} (no PyTorch, TensorFlow, or
	Keras) anywhere on the critical path, consistent with the CPU-only design.

	\begin{table}[t]
		\centering
		\caption{Software and library versions used for all builds and experiments.}
		\label{tab:setup-versions}
		\begin{tabular}{lll}
			\toprule
			Component & Version & Role \\
			\midrule
			Python            & 3.13.5   & runtime for the entire pipeline \\
			Neo4j             & 2026.05 (Community) & graph store (Dockerised) \\
			\code{neo4j} (driver) & 6.2.0 & Python$\to$Neo4j Bolt driver \\
			git               & 2.45.1   & reads per-commit file blobs / diffs \\
			javalang          & 0.13.0   & Java parser (AST/CFG/DFG/PDG/SEQ builders) \\
			NumPy             & 2.2.6    & arrays; the graph-method kernels \\
			SciPy             & 1.16.2   & sparse matrices; statistics (Mann--Whitney, ranks) \\
			scikit-learn      & 1.7.2    & logistic-regression heads, TF-IDF, RF/GB baselines \\
			pandas            & 2.3.2    & CSV / tabular handling \\
			networkx          & 3.5      & in-memory graph construction in the builders \\
			matplotlib        & 3.10.6   & figures \\
			nbformat / nbconvert & 5.10.4 / 7.17.1 & build and execute the result notebooks \\
			PyYAML            & 6.0.2    & per-project configuration (reproducibility package) \\
			\bottomrule
		\end{tabular}
	\end{table}

	\paragraph{A note on gensim / PyKEEN.} Earlier iterations experimented with
	\code{gensim} (node2vec/DeepWalk) and \code{PyKEEN} (TransE/DistMult). These are
	\textbf{not} used by the final methodology: the deployed DeepWalk embedding (DW) is a
	self-contained truncated-SVD of a PPMI matrix, and the deployed knowledge-graph
	embedding (KGE) is our own compact DistMult in NumPy (below). No external embedding
	library is on the final critical path.

	\subsection{Training hyper-parameters}
	\label{sec:setup-hparams}
	The checklist items \emph{optimiser}, \emph{learning rate}, \emph{batch size}, and
	\emph{number of epochs} are properties of gradient-descent neural training. Our final
	models are not neural, so most of these \textbf{do not apply}; the few genuine
	training knobs are given precisely.

	\begin{description}[leftmargin=1.4em]
		\item[Optimiser.] \emph{Mostly not applicable.} The linear heads on every graph
		method and on the $M$ and $G$ fusion channels are \textbf{balanced logistic
		regressions solved by L-BFGS} (scikit-learn, \code{solver="lbfgs"},
		\code{class\_weight="balanced"}, \code{max\_iter}$=1000$--$2000$); L-BFGS is a
		quasi-Newton batch solver, so there is no stochastic optimiser, no learning-rate
		schedule, and no mini-batching. The \emph{only} stochastic-gradient component in the
		whole system is the DistMult KGE, trained with plain SGD.
		\item[Learning rate.] \emph{Only the KGE has one:} the DistMult embedding uses a
		fixed SGD learning rate of $\eta=0.05$. All logistic-regression heads are L-BFGS and
		have no learning rate. DeepWalk (DW) is closed-form (SVD) and has none.
		\item[Batch size.] \emph{Not applicable in the neural sense.} There is no
		mini-batch training. The one gradient method (KGE) processes its past triples in
		internal chunks of $2{,}048$ for vectorised updates, but this is an implementation
		detail of a full-pass SGD, not a tuned batch-size hyper-parameter. The evaluation
		protocol does have a \emph{block} granularity ($200$ commits) at which models are
		refreshed, but that is a streaming-protocol constant (\secref{sec:setup-protocol}),
		not a training batch size.
		\item[Number of epochs.] \emph{Only the KGE has epochs:} the DistMult SGD runs
		$3$ epochs over the strictly-past triples at each refit; negative sampling uses
		$3$ corrupted tails per positive. No other model has an epoch count: the
		logistic heads run to L-BFGS convergence (capped at \code{max\_iter}), and DW/RN/PPR/LP
		are single-pass or fixed-iteration (below), not epoch-based.
	\end{description}

	\noindent The fixed-iteration counts of the non-embedding graph methods (which play the
	role ``epochs'' would in a neural model) are: \textbf{PPR} --- $40$ power-iterations of
	random-walk-with-restart, restart probability $\alpha=0.15$, run once per class seed set
	per block; \textbf{LP} --- $3$ propagation steps with clamping on known past labels;
	\textbf{RN} --- a single one-hop weighted vote (no iteration). These are held fixed
	across all projects and all graphs.

	\subsection{Model architecture}
	\label{sec:setup-arch}
	The checklist item \emph{hidden-layer dimensions} presupposes a multi-layer neural
	network. Our models have \textbf{no hidden layers} --- there is no multi-layer
	network, and specifically \textbf{no GNN} (a heterogeneous GNN is explicitly scoped as
	future work, not part of this study). The only learned latent dimensionalities are the
	two shallow embedding sizes already reported: \textbf{DeepWalk (DW) $=64$} and
	\textbf{DistMult KGE $=32$}. Beyond those:
	\begin{itemize}[nosep,leftmargin=1.4em]
		\item \textbf{RN, PPR, LP} produce a scalar risk score directly from the
		commit--hub incidence / transition matrix; they have no latent layer at all.
		\item \textbf{DW} embeds commits into $\mathbb{R}^{64}$ by truncated-SVD of the PPMI
		of the commit--hub co-occurrence (Levy--Goldberg form), then a logistic head; there
		is no hidden layer, only this one factorisation dimension.
		\item \textbf{KGE (DistMult)} embeds hubs and relations into $\mathbb{R}^{32}$ with a
		bilinear score; commits are folded in as the mean of their hubs' embeddings, then a
		logistic head. Again a single latent dimension, no hidden layers.
		\item \textbf{Text/structural feature widths} (not ``hidden'' layers, but the input
		dimensionalities worth recording): the structural TF--IDF and the CSTG text stream
		are hashed to $2^{18}$ features (a fixed, stateless hashing vectoriser); the CSTG
		typed-mass vector is $5$-dimensional; the JIT-metric channel is $12$-dimensional.
	\end{itemize}
	So the ``architecture'' of the deployed model is: five shallow graph-inference scorers
	(two of them with a $64$- or $32$-dim embedding) whose scalar outputs are combined by a
	single balanced logistic-regression stacker --- no hidden layers, no GNN.

	\subsection{Stopping and tuning}
	\label{sec:setup-stopping}
	\begin{description}[leftmargin=1.4em]
		\item[Early-stopping criteria.] \emph{Not used, and not applicable.} Because the
		models are non-neural, there is no validation-loss-based early stopping. The KGE
		SGD runs a fixed $3$ epochs; the logistic heads stop at L-BFGS convergence (or the
		\code{max\_iter} cap); the fixed-iteration methods (PPR $40$, LP $3$) run to their
		set iteration count. The streaming protocol itself has no ``stopping'': every
		post-warm-up commit is scored, and the run ends only when the commit stream ends.
		\item[Hyper-parameter tuning strategy.] We deliberately \textbf{do not tune the
		model hyper-parameters per project or per graph}, and this is a design choice, not
		an omission: the whole study is a \emph{controlled comparison}, so the protocol
		constants (warm-up $40\%$, block $200$, roll $800$, refit every $5$ blocks), the
		embedding sizes ($64$/$32$), the iteration counts (PPR $40$, LP $3$, KGE
		$3$ epochs / $\eta=0.05$ / $3$ negatives), and the classifier family (balanced L-BFGS
		logistic regression) are \textbf{held fixed across every representation, method,
		fusion, and project}. Fixing them is what makes cross-configuration and
		cross-project differences attributable to the factor under study rather than to
		per-cell tuning. The one thing that \emph{is} selected --- the deployed fusion
		formula --- is chosen not by a hyper-parameter search but by the transparent,
		reported combinatorial rule of \secref{sec:setup-fusion-selection} (evaluate all
		$2^5-1$ method combinations, take the smallest within a fixed tolerance of the best
		balanced score). The threshold at which hard decisions are read is likewise not a
		tuned constant but is re-derived \emph{online} from past-only data every $150$
		commits (\secref{sec:setup-protocol}). The default hyper-parameter values were set
		once, up front, from standard choices in the literature (e.g.\ $\alpha=0.15$ for
		PPR) and small sanity checks; there was no automated grid/random/Bayesian search
		over them, precisely to avoid over-fitting the evaluation.
	\end{description}

	\paragraph{Cross-project note.} Every constant above is identical across all subject
	projects (Groovy and the additional Apache projects processed with the reproducibility
	package); only the data changes. The single genuinely per-project quantity is the base
	commit that seeds the structural layers (\secref{sec:setup-data}), which is a property
	of each repository's history, not a tuned hyper-parameter.

\end{document} from the parent .tex.
% =====================================================================
\documentclass[11pt,a4paper]{article}
\usepackage[margin=2.4cm]{geometry}
\usepackage[T1]{fontenc}
\usepackage[utf8]{inputenc}
\usepackage{lmodern}
\usepackage{amsmath,amssymb}
\usepackage{booktabs}
\usepackage{array}
\usepackage{enumitem}
\usepackage{xcolor}
\usepackage[hidelinks]{hyperref}

% ---- macros used throughout (match the main document's conventions) ----
\providecommand{\code}[1]{\texttt{#1}}
\providecommand{\node}[1]{\textsf{#1}}
\providecommand{\edge}[1]{\textsc{#1}}

% ---- helper for safe section references (avoids any parsing issues with \S\ref) ----
\newcommand{\secref}[1]{\S\ref{#1}}

\begin{document}
	
	\section{Experimental Setup}
	\label{sec:experimental-setup}
	
	This section specifies, in full, the data, the knowledge-graph construction, the
	inference models, the online evaluation protocol, the metrics, the baselines, and
	the reproducibility conditions under which every result in this paper was produced.
	Two principles govern the whole design and are worth stating up front because they
	recur throughout. \textbf{First, everything is \emph{just-in-time (JIT) and fully
			online}:} the graph grows one commit at a time in true chronological order, and a
	commit is always predicted from a model and a graph state built \emph{only} from
	strictly earlier commits --- there is no random split, no look-ahead, and no
	statistic is ever fit on data that includes or postdates the commit under test.
	\textbf{Second, the study is designed as a controlled comparison:} when we ask which
	code representation or which inference method is best, exactly one factor is varied
	at a time while the protocol, the classifier family, the warm-up, and the block
	schedule are held fixed, so any measured difference is attributable to that factor
	and not to a confound.
	
	\subsection{Subject program, commit scope, and labels}
	\label{sec:setup-data}
	
	\paragraph{Project.} All experiments are conducted on \texttt{apache/groovy}, a
	mature, long-lived Java project (the Groovy programming language and its compiler,
	runtime, and tooling). We work from a full local \texttt{git} clone, so the raw
	material for every commit --- its message, its parent(s), and the exact byte content
	of every file before and after the change --- is available offline and deterministic.
	
	\paragraph{Labelled commit scope.} The study set is the \textbf{$8{,}059$ commits}
	of \texttt{apache/groovy} that appear in the \textbf{ApacheJIT} dataset
	(\texttt{data/apachejit/projects/apache\_groovy.csv}), spanning \textbf{2003--2019}.
	ApacheJIT provides, for each commit, a \emph{buggy/benign} label derived by an
	SZZ-style procedure (a commit is buggy if a later fix commit points back to it) and
	the twelve standard JIT change metrics. These $8{,}059$ commits are the only ones we
	treat as the labelled study set and the only ones grown into the AST/structural and
	semantic layers. The surrounding history is retained in the database as
	\node{Commit} nodes for context (the full clone has $16{,}887$ commits), but only the
	ApacheJIT commits are flagged \code{in\_jit=true} and are the units of prediction.
	
	\paragraph{Class balance and its consequences.} The labelled set is imbalanced:
	\textbf{$1{,}614$ buggy / $6{,}445$ benign}, i.e.\ a $\approx\!20\%$ positive rate
	overall. Crucially, the positive rate is \emph{not stationary}: it rises from
	$\approx\!6\%$ in 2003 to a peak of $\approx\!39\%$ in 2012 and then falls back to
	$\approx\!6\%$ by 2019 (measured directly from the labels). This pronounced
	\emph{concept drift} is the single most important property of the data for our
	purposes: it is why a random or single-split evaluation is inappropriate, why we
	adopt a time-aware online protocol, and why the tail of every online trajectory
	declines (the achievable precision/recall shrink as positives become rare). We
	return to this when defining the protocol (\secref{sec:setup-protocol}) and the
	metrics (\secref{sec:setup-metrics}).
	
	\paragraph{Base snapshot.} The structural layers are seeded once at a fixed
	\emph{base} commit --- \texttt{408b29851d7bbe4d343340832297e4be7e0c5578}, the first
	substantial Java commit --- whose $120$ \texttt{.java} files (${\approx}46$k AST
	nodes) constitute the initial code structure onto which all subsequent per-commit
	changes are laid. The base is a property of the \emph{code}, not of the labelled
	set; it merely fixes the starting state of the evolving graph.
	
	\subsection{Knowledge-graph construction and online growth}
	\label{sec:setup-kg}
	
	The knowledge graph is a \emph{heterogeneous information network} (HIN): typed nodes
	and typed edges in one store (Neo4j). It fuses three families of evidence that the
	defect-prediction literature usually treats separately --- \emph{process} evidence
	(who changed what, when, and how often), \emph{code-structure} evidence (the
	syntactic/semantic shape of the change and its surroundings), and \emph{textual}
	evidence (the natural-language and lexical content of the change). The design
	rationale for a graph (rather than a flat feature table) is that anchoring each
	commit's change to \emph{shared, persistent} code entities is what makes relational
	and neighbourhood reasoning across the entire history possible.
	
	\paragraph{Core (process / relational) layer.} The Core layer is the classical
	commit graph and is present in \emph{every} experimental configuration. Its node
	types are \node{Project}, \node{Branch}, \node{Developer}, \node{Commit},
	\node{File}, and \node{Issue}; its edges encode authorship
	(\node{Commit}$\xrightarrow{\text{\edge{AUTHORED\_BY}}}$\node{Developer}), the
	history spine (\node{Commit}$\xrightarrow{\text{\edge{PARENT\_OF}}}$\node{Commit}),
	file changes
	(\edge{ADDED}/\edge{MODIFIED}/\edge{DELETED}/\edge{RENAMED\_FROM}/\edge{RENAMED\_TO}
	between \node{Commit} and \node{File}), and issue links
	(\node{Commit}$\xrightarrow{\text{\edge{FIXES\_ISSUE}}}$\node{Issue}). Every
	\node{Commit} also carries the twelve ApacheJIT change metrics (size \code{la},
	\code{ld}; dispersion \code{nf}, \code{nd}, \code{ns}, \code{ent}; history and
	experience \code{age}, \code{ndev}, \code{nuc}, \code{aexp}, \code{arexp},
	\code{asexp}) and its \code{author\_ts} (the UNIX timestamp that defines online
	arrival order), \code{buggy} label, and \code{jit\_year}.
	
	\paragraph{Structural (code) layers as \emph{delta} graphs.} On top of the Core, we
	materialise the code structure of every changed \texttt{.java} file. The key modelling
	decision --- and a central contribution --- is that a commit is not attached to a
	per-commit snapshot blob; it is attached, by four typed \emph{delta} edges
	(\edge{ADDS}, \edge{REMOVES}, \edge{UPDATES}, \edge{MOVES}), directly to the specific
	code nodes it changed, laid on top of a single evolving structure per file. Concretely,
	for each tracked file the graph holds one \emph{current} structure; when a commit
	touches the file we (i) parse the before- and after-states with the \emph{same}
	deterministic builder, (ii) diff them, (iii) write the delta edges from the
	\node{Commit} to the affected nodes, and (iv) evolve the stored structure so the next
	commit diffs against the correct state. Node identity is threaded across versions by a
	deterministic, per-file parse-local-id map, so the ``after'' map of commit $N$ is
	byte-identical to the ``before'' map of the next commit that touches the file when the
	intervening source is unchanged. Removed nodes are \emph{retained} and flagged
	(\code{alive=false}, \code{removed\_by}), so the complete change history is queryable
	while the live structure always reflects the current code. This makes storage grow
	with the \emph{amount of change}, not with history-length $\times$ file-size, and it is
	what makes the whole graph online-extensible.
	
	We instantiate this delta-graph scheme for \textbf{five} structural representations of
	increasing coarseness, all built per changed \texttt{.java} file and attached at the
	\node{File} level:
	\begin{itemize}
		\item \textbf{AST} --- the abstract syntax tree; one node per syntactic construct
		(declaration, statement, expression, operator, literal, and each identifier/operator
		leaf). Finest granularity ($136$ node types).
		\item \textbf{CFG} --- the intra-procedural control-flow graph; one node per statement
		or control construct ($18$ node types).
		\item \textbf{DFG} --- the def--use data-flow graph over local variables; one node per
		definition/use site ($7$ node types).
		\item \textbf{PDG} --- the program-dependence graph; the CFG's statement vertices plus
		control-dependence and data-dependence \emph{edges} (no new nodes; $18$ node types).
		We label it PDG rather than CPG because we deliberately keep dependence at statement
		granularity and do \emph{not} embed the AST subtree under each statement; a canonical
		code property graph would contain the AST as a subgraph and is out of scope.
		\item \textbf{Token-seq} --- a deliberately weak flat statement sequence (statement
		nodes chained by \edge{NEXT}, no control or data structure); a lower-bound rival used
		only in the representation comparison.
	\end{itemize}
	Table~\ref{tab:setup-layers} reports the materialised scale of each layer. The
	${\approx}10\times$ gap between AST ($3.70$\,M nodes) and the others ($0.26$--$0.47$\,M)
	is a pure \emph{granularity} effect: a statement such as \code{int x = foo(a,b)+1;} is
	one node in CFG/PDG but a dozen-plus nodes in the AST. CFG and PDG have identical node
	counts because the PDG adds only dependence edges to the CFG's vertex set.
	
	\begin{table}[t]
		\centering
		\caption{Materialised scale of each structural layer and its per-commit change-token
			stream (Groovy, $8{,}059$ commits). ``Token types'' counts the distinct
			\code{\{edge\}:\{node-type\}} change-tokens; ``Commits w/ tok.'' is the number of
			labelled commits that emit at least one token in that layer.}
		\label{tab:setup-layers}
		\begin{tabular}{lrrrrr}
			\toprule
			Layer & Nodes & Delta-edges & Node types & Token types & Commits w/ tok. \\
			\midrule
			AST       & 3{,}699{,}439 & 3{,}614{,}922 & 136 & 314 & 7{,}523 \\
			CFG       &   372{,}488 &   408{,}301 & 18 & 53 & 6{,}194 \\
			DFG       &   469{,}662 &   648{,}553 &  7 & 27 & 6{,}066 \\
			PDG       &   372{,}488 &   411{,}869 & 18 & 63 & 6{,}261 \\
			Token-seq &   257{,}708 &   293{,}287 & 16 & 47 & 6{,}011 \\
			\bottomrule
		\end{tabular}
	\end{table}
	
	\paragraph{The differ, and why matching is fair across layers.} The AST layer is
	diffed as a \emph{tree} by bottom-up subtree hashing (identical subtrees in the before
	and after trees are matched greedily, preferring equal source positions; unmatched
	before-nodes are deletions, unmatched after-nodes insertions, matched-but-relabelled
	nodes updates, reparented nodes moves). The four non-tree layers (CFG/DFG/PDG/Token-seq)
	are general typed graphs and are diffed by a single \emph{node-identity-lite} matcher
	applied \emph{identically} to each: matching is scoped within a method (\code{class::name},
	stable across commits) and keyed by \emph{ordinal} (the $k$-th node of a given type in a
	method maps to the $k$-th such node in the other version), which is robust to line
	shifts. Using one shared differ for the alternatives is the fairness guarantee for the
	representation comparison: only the subgraph varies, never the difference model. We
	disclose that the AST uses a tree differ while the others use the ordinal matcher --- a
	deliberate specialisation (trees vs.\ general graphs), not a hidden confound.
	
	\paragraph{Semantic-text layer (CSTG).} The third first-class layer is the
	\textbf{Commit Semantic-Text Graph (CSTG)}, which models the \emph{textual}
	description of a change (commit message $+$ diff text) as a graph grounded in code.
	Its node types are \node{Term} (a typed vocabulary term, with \code{kind} $\in$
	\{\emph{code}-identifier, \emph{bug}-lexicon, \emph{action}, \emph{error}-type,
	\emph{nl}\}) and \node{Intent} (the change-intent class: fix, feat, refactor, \dots);
	its edges are
	\node{Commit}$\xrightarrow{\text{\edge{MENTIONS}\,(w=\text{TW-IDF})}}$\node{Term} (top
	terms per commit, weighted by TextRank centrality in the commit's word-graph $\times$
	IDF), \node{Term}$\xrightarrow{\text{\edge{COOCCURS}\,(\text{npmi})}}$\node{Term} (the
	global normalised-pointwise-mutual-information term graph),
	\node{Term}$\xrightarrow{\text{\edge{GROUNDS\_IN}}}$\node{ASTNode} (a code term linked
	to the alive identifier leaf it names), and
	\node{Commit}$\xrightarrow{\text{\edge{HAS\_INTENT}}}$\node{Intent}. On Groovy the CSTG
	comprises $20{,}377$ \node{Term} nodes, $188{,}435$ \edge{MENTIONS}, $120{,}340$
	\edge{COOCCURS}, $19{,}315$ \edge{GROUNDS\_IN}, and $8{,}059$ \edge{HAS\_INTENT}.
	Theoretically the layer rests on graph-of-words / TW-IDF (Rousseau \& Vazirgiannis),
	the NPMI corpus term graph (cf.\ TextGCN), defect-semantic typing, and term-risk
	propagation over the NPMI graph. Every CSTG statistic that enters prediction is grown
	online, past-only (see \secref{sec:setup-protocol}); the static term graph and
	term-risk values written to Neo4j are a queryable/visualisation snapshot and are
	\emph{not} read by the online predictor.
	
	\paragraph{Data-quality guarantees.} The finalised graph passes a standing integrity
	audit: no file has more than one structure root, removed nodes always carry a
	\edge{REMOVES} edge, and the online-evolved structure reproduces a fresh
	ground-truth parse of the current file state with precision $=$ recall $=1.0$ on the
	validated sample. A large \emph{ghost-AST} inflation ($\approx\!1.58$\,M nodes left by
	package-restructuring renames outside the labelled set) was identified and corrected,
	so the reported node counts reflect the clean final state. Known, documented
	limitations that a reviewer should be aware of: labels are SZZ-derived and carry the
	usual mining noise; non-\texttt{.java} files are not modelled at the code level; and
	a git rename is currently treated as an add-of-new-path in the delta layer.
	
	\subsection{Inference models}
	\label{sec:setup-models}
	
	The target hardware has \textbf{no GPU}; every method is therefore deliberately
	CPU-only and \textbf{GNN-free}, both to be deployable and to constitute strong,
	honest baselines that any future GNN must beat. The final methodology fixes
	\textbf{five canonical graph-inference paradigms}, chosen because they map one-to-one
	onto the recognised families of graph and knowledge-graph inference, each reads a
	\emph{different, complementary} facet of the structure, and --- crucially --- each is a
	graph \emph{algorithm} that runs unchanged on \emph{any} graph in the family, so all
	five are valid rows across every column of the representation study:
	
	\begin{description}
		\item[RN --- Relational neighbour (wvRN).] Collective classification: a commit's
		score is the label-weighted hub-overlap with past commits (one hop through the shared
		file/developer/term hubs). Complexity $O(\mathrm{nnz})$ of a sparse block matmul.
		\item[PPR --- Personalised PageRank (random walk with restart).] The bipartite
		commit--hub graph is seeded from past buggy vs.\ benign commits and a commit is scored
		by its relative proximity to each class. Complexity $O(\text{iters}\cdot|E|)$ per
		block, run once per class seed set.
		\item[LP --- Label propagation.] Spreading activation of past labels through the
		hubs a few steps, clamped on known past labels. Complexity $O(\text{iters}\cdot|E|)$.
		\item[DW --- DeepWalk-style embedding.] A random-walk / matrix-factorisation node
		embedding (PPMI of the commit--hub co-occurrence, truncated-SVD, the Levy--Goldberg
		form) with a logistic head; marginals from past rows only.
		\item[KGE --- Knowledge-graph embedding (DistMult).] A compact bilinear embedding
		over typed $(\text{commit},\text{rel},\text{hub})$ triples restricted to past
		commits, with commits folded in from their hubs' learned embeddings and a logistic
		head.
	\end{description}
	
	The embedding methods (DW, KGE) are refit on the expanding past window every
	\textbf{$5$ blocks}; RN, PPR, and LP are recomputed each block from the strictly-past
	state. Embedding dimensions are $64$ (DW) and $32$ (KGE). All linear heads are
	balanced logistic regressions.
	
	\paragraph{Two non-graph channels used only in fusion.} Two additional signals are
	\emph{not} admissible as rows of the representation study (they are not graph
	algorithms) but are evaluated as fusion channels:
	\begin{itemize}
		\item \textbf{$M$ --- JIT metrics:} a prequential logistic model over the twelve
		ApacheJIT change metrics; the classical JIT-defect predictor and our metric-only
		reference channel.
		\item \textbf{$G$ --- CSTG feature channel:} a prequential logistic model over the
		CSTG's per-commit features --- the graph-of-words (TextRank-weighted, add/remove
		polarity-tagged) term stream, the NPMI-propagated past-only term-risk prior, a
		$5$-dimensional typed-mass vector over the defect lexicon, and the intent/consistency
		features --- enriched with an online recency / bug-cache block (per-file and
		per-developer time-since-last-buggy touch, recently-buggy touched-file count, re-fix
		flag, commit velocity), each computed from strictly-past state.
	\end{itemize}
	The commit-\emph{message} text channel ($X$) that appeared in earlier iterations is
	\textbf{removed entirely} from the final methodology (it is a strict subset of $G$ and
	adds nothing over it); it is disabled throughout and reported only as a historical
	step.
	
	\paragraph{Fusion.} Method and channel scores are combined by \emph{prequential
		logistic-regression stacking}: at each block a balanced logistic regression is fit on
	the strictly-past predicted probabilities of the constituent models and applied to the
	current block, refit on the expanding window. The \emph{deployed} model is selected by
	an explicit load-vs-performance rule (\secref{sec:setup-fusion-selection}) and is the
	purely graph-native \textbf{$F{+}G = \text{RN}{+}\text{PPR}{+}\text{CSTG}$}, i.e.\ the
	two-method graph fusion $F=\text{RN}{+}\text{PPR}$ augmented with the CSTG semantic
	channel $G$, \emph{without} the non-graph JIT metrics $M$.
	
	\subsection{The two research questions and the graph family}
	\label{sec:setup-rqs}
	
	The experiments answer two questions with one design:
	\textbf{(RQ1, which representation?)} does a richer code/semantic structure improve
	commit-level prediction, and which structural representation is best?
	\textbf{(RQ2, which inference method?)} among typical graph-inference paradigms, which
	best exploits the graph?
	
	Both are answered by a single family of graphs of increasing richness, all sharing the
	Core relational layer, evaluated with the fixed five methods:
	\emph{Core}; \emph{Core$+$AST}; \emph{Core$+$CFG}; \emph{Core$+$DFG};
	\emph{Core$+$PDG}; and the deployed final \emph{Core$+$AST$+$CSTG}. Reading a row
	(a method across graphs) answers RQ1 for that method; reading a column (all methods on
	one graph) answers RQ2 for that graph. A structural layer influences prediction through
	exactly one channel --- the per-commit change-tokens \code{\{edge\}:\{type\}} that
	become typed hubs the walks and embeddings traverse --- so swapping the representation
	means swapping the hub stream while everything else is held fixed. On the final graph,
	the hubs additionally include CSTG \node{Term} nodes via \edge{MENTIONS}, so the same
	five methods traverse the semantic signal too. (Token-seq is included in the
	representation \emph{ablation} but is not part of the final six-graph family, being a
	deliberately weak rival.)
	
	\subsection{Online evaluation protocol}
	\label{sec:setup-protocol}
	
	\paragraph{Prequential, predict-then-grow.} Commits are streamed in chronological
	order (\code{author\_ts}). The protocol is \emph{prequential}: for each commit we
	(i)~predict its label from a model and graph state built only from strictly earlier
	commits, (ii)~reveal the true label, (iii)~grow the graph state (add the commit's
	delta edges, term hubs, recency counters, co-occurrence counts), and (iv)~periodically
	re-learn on the expanded window. This exactly reproduces the information available at
	deployment time and structurally precludes look-ahead leakage.
	
	\paragraph{Warm-up, blocks, and refit schedule.} The first \textbf{$40\%$} of the
	stream ($W = \lfloor 0.40\times 8{,}059\rfloor = 3{,}223$ commits) is a warm-up used
	only to initialise models and is \emph{not} scored; the remaining $4{,}836$ commits are
	the evaluation stream. Predictions are made per commit but models are refreshed on a
	\textbf{block} granularity of $200$ commits (the expanding-window retraining that lets
	the models adapt to concept drift). This warm-up/block schedule, the classifier family,
	and all hyper-parameters are held \emph{identical} across every representation, method,
	and fusion so that the only independent variable in each comparison is the one under
	study.
	
	\paragraph{Leakage controls (stated explicitly).} All scalers, vectorisers, IDF
	weights, NPMI graphs, term-risk priors, embedding bases, and stacking regressions are
	fit on past data only; PPR class seeds never include the commit under test; the
	graph-of-words weighting and recency lookups are commit-local; and the
	\code{HashingVectorizer} used for the CSTG text stream is stateless (a fixed hash with
	no learned vocabulary), so it cannot transport information from a future commit into an
	earlier row. The single-split controlled comparison used against the external text
	baseline (\secref{sec:setup-baselines}) is clearly separated from the deployed online
	path and never mixed with it.
	
	\paragraph{Online operating point.} Because the task is imbalanced and drifting, the
	hard-decision threshold is itself tuned \emph{online}: it is re-tuned on the past-only
	predictions every $150$ commits to maximise buggy-F1 on already-seen data, giving a
	leakage-free deployment operating point at which all threshold-based metrics are read.
	Rolling trajectories are computed over a trailing window of $\text{ROLL}=800$ commits
	(a visualisation smoothing only --- it never enters the model or the reported scalars,
	which are computed once over the entire evaluation stream).
	
	\subsection{Evaluation metrics}
	\label{sec:setup-metrics}
	
	Because the positive class is rare and drifting, no single metric suffices; we report a
	\textbf{seven-metric} suite plus two effort-aware measures, and treat the
	threshold-free and class-balanced metrics as primary.
	
	\begin{itemize}
		\item \textbf{Threshold-free ranking:} \emph{ROC-AUC} (a.k.a.\ AUC) and, where
		reported, \emph{PR-AUC} (area under precision--recall; its random baseline equals the
		positive rate, which is why absolute PR-AUC falls in the low-defect tail even for a
		good model).
		\item \textbf{Operating-point metrics at the online-tuned threshold:} \emph{Precision},
		\emph{Recall} (buggy-class sensitivity), \emph{Buggy-F1} (F1 of the positive class),
		\emph{Macro-F1} (unweighted mean of both classes' F1 --- robust to imbalance),
		\emph{G-Mean} ($\sqrt{\text{TPR}\cdot\text{TNR}}$), and \emph{MCC} (Matthews
		correlation --- the single most reliable summary under imbalance).
		\item \textbf{Effort-aware (deployment triage) metrics:} normalised
		$P_{\mathrm{opt}}$ (Kamei et al.: how close the model's inspection order is to the
		ideal buggy-first, least-effort-first order, normalised so $0.5$ is random and $1.0$
		optimal) and \emph{ACC@20\%LOC} (recall achieved when only the top $20\%$ of churned
		lines are inspected). Effort is measured as \code{la}$+$\code{ld}.
	\end{itemize}
	The two \emph{class-balanced} metrics --- Macro-F1 and G-Mean --- are the ones we
	prioritise for model selection, precisely because they do not let a model win by
	exploiting the majority class in the low-defect regime. We emphasise that, given the
	concept drift, the meaningful reading of an online trajectory is \emph{relative} (which
	method/graph stays highest) rather than the absolute end-of-stream value, since the
	falling positive rate depresses every method's tail together.
	
	\subsection{Fusion selection rule}
	\label{sec:setup-fusion-selection}
	
	The deployed fusion is not hand-picked. We evaluate \emph{all} $2^5-1=31$ combinations
	of the five methods on the final graph, score each on the full seven-metric suite, and
	select the smallest combination whose balanced score
	$(\text{Macro-F1}+\text{G-Mean})/2$ is within a fixed tolerance ($0.005$) of the best ---
	a transparent trade-off of \emph{load} (fewer methods, cheaper online) against
	\emph{performance}. This selects $F=\text{RN}{+}\text{PPR}$. We then test, on top of
	$F$, whether the explicit CSTG channel $G$ and/or the non-graph JIT metrics $M$ add
	anything: $G$ helps markedly and $M$ does not (adding $M$ alone hurts), so the deployed
	model is $F{+}G=\text{RN}{+}\text{PPR}{+}\text{CSTG}$. Reporting the full combinatorial
	ablation, and choosing by an explicit rule, is what makes the deployed model
	defensible rather than cherry-picked.
	
	\subsection{Baselines and comparison protocols}
	\label{sec:setup-baselines}
	
	We compare against baselines at two levels.
	\textbf{(i) Within-project JIT baselines:} logistic regression, random forest, and
	gradient boosting over the twelve ApacheJIT metrics --- the standard JIT-defect
	predictor --- evaluated under the identical online protocol, plus a naive prior-rate
	predictor. \textbf{(ii) A strong external text baseline:} a $768$-dimensional
	diff-text TF-IDF model with a random forest, the strongest published lexical baseline.
	Because that baseline was reported on a single chronological split (train first $70\%$,
	test last $15\%$; a hard, low-defect tail with $\approx\!7.5\%$ positives), we
	reproduce that \emph{exact} split and test set for a like-for-like comparison, using a
	matched classifier so any difference comes from the features and not the learner; these
	single-split numbers are reported separately from, and never conflated with, the
	fully-online results. As a sanity check our JIT-metrics model reproduces the baseline
	logistic regression on that split to three decimals, confirming the protocols are
	aligned. For the CSTG internal ablation specifically, we additionally report a
	\emph{same-learner} (matched random forest) plain-TF-IDF row so the ``does the CSTG
	representation beat a bag-of-words'' comparison is under one learner rather than mixing
	linear and non-linear models.
	
	\subsection{Scalability, complexity, and statistical-significance setup}
	\label{sec:setup-scalability}
	
	Beyond accuracy, the paper characterises the system's cost and the statistical
	reliability of its comparisons. All of this is computed \emph{without rebuilding} the
	knowledge graph: final sizes and growth come from read-only queries against the built
	graph and from attributing delta edges to their source commits in online order;
	prediction latency is measured by replaying the cached graph family in memory; and
	per-commit build/modify cost is obtained by an instrumented replay of the real
	parse-and-diff work on a stratified sample of commits, cross-checked against an
	analytical operation-count proxy.
	
	\begin{itemize}
		\item \textbf{Statistical profile.} Per-layer node/edge counts, type and token
		vocabularies, per-file size distributions (with Gini and skew), and a discrete
		power-law fit of the per-commit change-size distribution (heavy-tailed, exponent
		$2<\alpha<3$ on every layer).
		\item \textbf{Per-commit growth and modification.} Cumulative nodes/edges over the
		commit stream and over calendar time (overlaid on the drifting bug rate), the
		per-commit change-size distribution, and the buggy-vs-benign change-size separation
		--- which is large and highly significant (Mann--Whitney $p\ll 10^{-100}$ on every
		layer; buggy commits make on the order of $2$--$3\times$ larger structural changes),
		the basic reason the graph is predictive.
		\item \textbf{Time and space complexity.} Build/modify cost is shown to be near-linear
		in file size ($R^2\approx0.99$), i.e.\ per-commit work is $O(\text{change})$ and space
		is $O(\text{alive nodes}+\text{retained history})$; prediction latency scales with
		graph density, and the deployed graph-native fusion predicts at sub-millisecond per
		commit (thousands of commits/second), so the whole pipeline is comfortably real-time
		and GPU-free. Each per-commit method's cost is annotated with its known asymptotic
		complexity and validated against the measured block-time-vs-graph-size growth.
		\item \textbf{Statistical significance.} Headline comparisons are tested on paired
		per-commit out-of-sample scores: DeLong's test for ROC-AUC differences (Holm-corrected
		within each family of comparisons), paired bootstrap tests and $95\%$ confidence
		intervals for PR-AUC/MCC/G-Mean, and Cliff's~$\delta$ effect sizes. This lets us state
		which accuracy gaps are significant (e.g.\ graph richness significantly helps the
		informative methods; PPR is significantly the best method on the final graph; the
		semantic channel $F{+}G$ significantly improves every headline metric over $F$) and
		which are within noise.
	\end{itemize}
	
	\subsection{Implementation and reproducibility}
	\label{sec:setup-impl}
	
	The graph store is Neo4j (final on-disk footprint $\approx\!3.3$\,GB; working footprint
	$\approx\!5.6$\,GB with transaction logs). Java is parsed with a single deterministic
	parser (\texttt{javalang}) for the whole pipeline, in killable subprocesses with a
	watchdog so a pathological file cannot stall the build; a persistent parser worker and
	an in-memory before-graph cache keep the online build tractable, and a range index on
	each node label's \code{id} keeps per-commit cost constant as the graph grows (an
	absent index was, at one point, the dominant cost --- turning an $O(\text{change})$
	build into an accidentally quadratic one). All inference code is CPU-only Python
	(scikit-learn linear models, custom sparse graph routines, gensim/PyKEEN for the
	embeddings where used); no GNN and no GPU are required. Feature streams and per-method
	scores are cached to disk so every experiment, table, figure, and statistical test can
	be regenerated without a live database or a graph rebuild, and companion Jupyter
	notebooks reproduce all figures from those caches. Randomised components (embedding
	initialisation, bootstrap resampling) use fixed seeds. Determinism of node identity
	across versions (the same source always yields the same parse-local ids) is what makes
	the online build itself reproducible.
	
	\paragraph{Summary of the final configuration.} The deployed predictor is a
	fully-online, GPU-free, multimodal commit knowledge graph on \texttt{apache/groovy}
	($8{,}059$ ApacheJIT commits, $\approx\!20\%$ buggy, 2003--2019): a Core process layer,
	an AST delta layer, and the CSTG semantic-text layer, over which the fusion
	$F{+}G=\text{RN}{+}\text{PPR}{+}\text{CSTG}$ predicts each commit prequentially from
	strictly-past state, warm-up $40\%$ ($3{,}223$ commits), block $200$, evaluated on the
	remaining $4{,}836$ commits with a seven-metric plus effort-aware suite and paired
	significance testing.

	% =====================================================================
	\section{Environment, Software Versions, and Model-Configuration Checklist}
	\label{sec:setup-checklist}
	% =====================================================================
	For completeness and reproducibility, this section answers a standard checklist of
	setup details --- operating system and hardware, software and library versions,
	training hyper-parameters, model architecture, and stopping/tuning. Several checklist
	items are framed for \emph{deep-learning} pipelines (optimiser, learning rate, batch
	size, epochs, hidden-layer dimensions, early stopping). Our final methodology is
	deliberately \textbf{GPU-free and GNN-free}: the inference models are classical
	graph algorithms (random-walk / label-propagation / matrix-factorisation) with
	\emph{linear} (logistic-regression) heads, plus one compact bilinear knowledge-graph
	embedding trained by a short custom SGD. Where a checklist item has no analogue in
	this design, we say so explicitly and explain why, rather than inventing a value.

	\subsection{Operating system and hardware}
	\label{sec:setup-hw}
	All building and all experiments were run on a single commodity laptop, with no GPU
	and no cluster:
	\begin{itemize}[nosep,leftmargin=1.4em]
		\item \textbf{Operating system:} Microsoft Windows~11 (64-bit, build $26{,}200$).
		The pipeline is OS-agnostic (pure Python plus a Dockerised Neo4j) and also runs on
		Linux/macOS; the Neo4j administration helpers assume a Docker-hosted server.
		\item \textbf{CPU:} Intel Core i7-13620H (13th generation), $10$ physical cores /
		$16$ logical threads, $2.4$\,GHz base clock. \textbf{No GPU is used at any stage}
		(build, inference, or analysis).
		\item \textbf{Memory (RAM):} $16$\,GB. The whole pipeline is designed to run within
		this budget: the online build stores only one evolving structure per file (not a
		snapshot per commit), and the inference/analysis code streams and caches to disk, so
		no stage requires loading the full graph into memory at once.
		\item \textbf{Graph store:} Neo4j runs in a Docker container; the built Groovy graph
		occupies $\approx\!3.3$\,GB on disk ($\approx\!5.6$\,GB with transaction logs).
	\end{itemize}

	\subsection{Software and library versions}
	\label{sec:setup-versions}
	The exact versions used are listed in Table~\ref{tab:setup-versions}. The final
	pipeline depends only on the mainstream scientific-Python stack, a Java parser, and a
	Neo4j driver; there is \textbf{no deep-learning framework} (no PyTorch, TensorFlow, or
	Keras) anywhere on the critical path, consistent with the CPU-only design.

	\begin{table}[t]
		\centering
		\caption{Software and library versions used for all builds and experiments.}
		\label{tab:setup-versions}
		\begin{tabular}{lll}
			\toprule
			Component & Version & Role \\
			\midrule
			Python            & 3.13.5   & runtime for the entire pipeline \\
			Neo4j             & 2026.05 (Community) & graph store (Dockerised) \\
			\code{neo4j} (driver) & 6.2.0 & Python$\to$Neo4j Bolt driver \\
			git               & 2.45.1   & reads per-commit file blobs / diffs \\
			javalang          & 0.13.0   & Java parser (AST/CFG/DFG/PDG/SEQ builders) \\
			NumPy             & 2.2.6    & arrays; the graph-method kernels \\
			SciPy             & 1.16.2   & sparse matrices; statistics (Mann--Whitney, ranks) \\
			scikit-learn      & 1.7.2    & logistic-regression heads, TF-IDF, RF/GB baselines \\
			pandas            & 2.3.2    & CSV / tabular handling \\
			networkx          & 3.5      & in-memory graph construction in the builders \\
			matplotlib        & 3.10.6   & figures \\
			nbformat / nbconvert & 5.10.4 / 7.17.1 & build and execute the result notebooks \\
			PyYAML            & 6.0.2    & per-project configuration (reproducibility package) \\
			\bottomrule
		\end{tabular}
	\end{table}

	\paragraph{A note on gensim / PyKEEN.} Earlier iterations experimented with
	\code{gensim} (node2vec/DeepWalk) and \code{PyKEEN} (TransE/DistMult). These are
	\textbf{not} used by the final methodology: the deployed DeepWalk embedding (DW) is a
	self-contained truncated-SVD of a PPMI matrix, and the deployed knowledge-graph
	embedding (KGE) is our own compact DistMult in NumPy (below). No external embedding
	library is on the final critical path.

	\subsection{Training hyper-parameters}
	\label{sec:setup-hparams}
	The checklist items \emph{optimiser}, \emph{learning rate}, \emph{batch size}, and
	\emph{number of epochs} are properties of gradient-descent neural training. Our final
	models are not neural, so most of these \textbf{do not apply}; the few genuine
	training knobs are given precisely.

	\begin{description}[leftmargin=1.4em]
		\item[Optimiser.] \emph{Mostly not applicable.} The linear heads on every graph
		method and on the $M$ and $G$ fusion channels are \textbf{balanced logistic
		regressions solved by L-BFGS} (scikit-learn, \code{solver="lbfgs"},
		\code{class\_weight="balanced"}, \code{max\_iter}$=1000$--$2000$); L-BFGS is a
		quasi-Newton batch solver, so there is no stochastic optimiser, no learning-rate
		schedule, and no mini-batching. The \emph{only} stochastic-gradient component in the
		whole system is the DistMult KGE, trained with plain SGD.
		\item[Learning rate.] \emph{Only the KGE has one:} the DistMult embedding uses a
		fixed SGD learning rate of $\eta=0.05$. All logistic-regression heads are L-BFGS and
		have no learning rate. DeepWalk (DW) is closed-form (SVD) and has none.
		\item[Batch size.] \emph{Not applicable in the neural sense.} There is no
		mini-batch training. The one gradient method (KGE) processes its past triples in
		internal chunks of $2{,}048$ for vectorised updates, but this is an implementation
		detail of a full-pass SGD, not a tuned batch-size hyper-parameter. The evaluation
		protocol does have a \emph{block} granularity ($200$ commits) at which models are
		refreshed, but that is a streaming-protocol constant (\secref{sec:setup-protocol}),
		not a training batch size.
		\item[Number of epochs.] \emph{Only the KGE has epochs:} the DistMult SGD runs
		$3$ epochs over the strictly-past triples at each refit; negative sampling uses
		$3$ corrupted tails per positive. No other model has an epoch count: the
		logistic heads run to L-BFGS convergence (capped at \code{max\_iter}), and DW/RN/PPR/LP
		are single-pass or fixed-iteration (below), not epoch-based.
	\end{description}

	\noindent The fixed-iteration counts of the non-embedding graph methods (which play the
	role ``epochs'' would in a neural model) are: \textbf{PPR} --- $40$ power-iterations of
	random-walk-with-restart, restart probability $\alpha=0.15$, run once per class seed set
	per block; \textbf{LP} --- $3$ propagation steps with clamping on known past labels;
	\textbf{RN} --- a single one-hop weighted vote (no iteration). These are held fixed
	across all projects and all graphs.

	\subsection{Model architecture}
	\label{sec:setup-arch}
	The checklist item \emph{hidden-layer dimensions} presupposes a multi-layer neural
	network. Our models have \textbf{no hidden layers} --- there is no multi-layer
	network, and specifically \textbf{no GNN} (a heterogeneous GNN is explicitly scoped as
	future work, not part of this study). The only learned latent dimensionalities are the
	two shallow embedding sizes already reported: \textbf{DeepWalk (DW) $=64$} and
	\textbf{DistMult KGE $=32$}. Beyond those:
	\begin{itemize}[nosep,leftmargin=1.4em]
		\item \textbf{RN, PPR, LP} produce a scalar risk score directly from the
		commit--hub incidence / transition matrix; they have no latent layer at all.
		\item \textbf{DW} embeds commits into $\mathbb{R}^{64}$ by truncated-SVD of the PPMI
		of the commit--hub co-occurrence (Levy--Goldberg form), then a logistic head; there
		is no hidden layer, only this one factorisation dimension.
		\item \textbf{KGE (DistMult)} embeds hubs and relations into $\mathbb{R}^{32}$ with a
		bilinear score; commits are folded in as the mean of their hubs' embeddings, then a
		logistic head. Again a single latent dimension, no hidden layers.
		\item \textbf{Text/structural feature widths} (not ``hidden'' layers, but the input
		dimensionalities worth recording): the structural TF--IDF and the CSTG text stream
		are hashed to $2^{18}$ features (a fixed, stateless hashing vectoriser); the CSTG
		typed-mass vector is $5$-dimensional; the JIT-metric channel is $12$-dimensional.
	\end{itemize}
	So the ``architecture'' of the deployed model is: five shallow graph-inference scorers
	(two of them with a $64$- or $32$-dim embedding) whose scalar outputs are combined by a
	single balanced logistic-regression stacker --- no hidden layers, no GNN.

	\subsection{Stopping and tuning}
	\label{sec:setup-stopping}
	\begin{description}[leftmargin=1.4em]
		\item[Early-stopping criteria.] \emph{Not used, and not applicable.} Because the
		models are non-neural, there is no validation-loss-based early stopping. The KGE
		SGD runs a fixed $3$ epochs; the logistic heads stop at L-BFGS convergence (or the
		\code{max\_iter} cap); the fixed-iteration methods (PPR $40$, LP $3$) run to their
		set iteration count. The streaming protocol itself has no ``stopping'': every
		post-warm-up commit is scored, and the run ends only when the commit stream ends.
		\item[Hyper-parameter tuning strategy.] We deliberately \textbf{do not tune the
		model hyper-parameters per project or per graph}, and this is a design choice, not
		an omission: the whole study is a \emph{controlled comparison}, so the protocol
		constants (warm-up $40\%$, block $200$, roll $800$, refit every $5$ blocks), the
		embedding sizes ($64$/$32$), the iteration counts (PPR $40$, LP $3$, KGE
		$3$ epochs / $\eta=0.05$ / $3$ negatives), and the classifier family (balanced L-BFGS
		logistic regression) are \textbf{held fixed across every representation, method,
		fusion, and project}. Fixing them is what makes cross-configuration and
		cross-project differences attributable to the factor under study rather than to
		per-cell tuning. The one thing that \emph{is} selected --- the deployed fusion
		formula --- is chosen not by a hyper-parameter search but by the transparent,
		reported combinatorial rule of \secref{sec:setup-fusion-selection} (evaluate all
		$2^5-1$ method combinations, take the smallest within a fixed tolerance of the best
		balanced score). The threshold at which hard decisions are read is likewise not a
		tuned constant but is re-derived \emph{online} from past-only data every $150$
		commits (\secref{sec:setup-protocol}). The default hyper-parameter values were set
		once, up front, from standard choices in the literature (e.g.\ $\alpha=0.15$ for
		PPR) and small sanity checks; there was no automated grid/random/Bayesian search
		over them, precisely to avoid over-fitting the evaluation.
	\end{description}

	\paragraph{Cross-project note.} Every constant above is identical across all subject
	projects (Groovy and the additional Apache projects processed with the reproducibility
	package); only the data changes. The single genuinely per-project quantity is the base
	commit that seeds the structural layers (\secref{sec:setup-data}), which is a property
	of each repository's history, not a tuned hyper-parameter.

\end{document}